\chapter{Lo que falta pedir}

\label{cap:pedidos}
Este apéndice reúne, ordenado por destinatario y no por capítulo, todo lo que el libro
identificó como faltante. No es una lista de deseos: cada entrada nombra un documento
existente, con su organismo y —cuando se lo conoce— su número de expediente. Está pensado
para que quien continúe la investigación pueda entrar a una mesa de entradas con el pedido
ya escrito.

Entre paréntesis, el capítulo donde la falta se explica. Cuando un pedido tiene más de un destinatario, figura en el principal. Se unificaron los pedidos que se repetían ante un mismo destinatario; cuando un documento puede pedirse por más de una vía ---la matrícula 4366, por ejemplo, en Inmuebles y en el expediente hídrico--- se lo mantiene en cada una.

\section{Dirección General de Inmuebles}

\begin{itemize}[leftmargin=*]
\item \textbf{antecedentes de la finca «El Angosto»} deslindada en 1921 a pedido de Antonia P.\ de Peralta, y su correspondencia con la parcela que hoy lleva ese nombre en el campo «finca» del catastro: sus cuatro linderos de 1921 ---Daniel Linares\index{Linares, Daniel}, Julio Mercado, Susana Aguilera y un Augusto Reyes--- repiten tres de los que el remate del Angosto de Arrieta declaraba en 1909, y el cotejo de las dos superficies es lo único que puede decir si son la misma tierra \textit{(caps.\ \ref{cap:fincas} y \ref{cap:siglo})}

\item \textbf{matrícula e informe de dominio del loteo cerrado Solanas de La Caldera} ---cuyo ingreso sobre la Ruta Nacional 9, a la altura del kilómetro 1618, cedió el 4 de enero de 2026---: titulares registrales, porcentajes y antecedente dominial. Es el documento que confirma o descarta la copropiedad que la prensa atribuye a la actual Secretaria de Tierras y Bienes del Estado, y ninguna afirmación sobre ella debe publicarse sin él \textit{(caps.\ \ref{cap:loteo} y \ref{cap:opacidad})}
\item \textbf{matrícula y antecedente dominial de cualquiera de las parcelas en cinta del valle al sur del casco urbano} ---las que la lámina \ref{lam:cintas} muestra---, para establecer \textbf{si provienen de una partición agrícola antigua o de una subdivisión reciente}. De esa respuesta depende que el apartado correspondiente del capítulo \ref{cap:loteo} se confirme o se elimine \textit{(cap.\ \ref{cap:loteo})}
\item \textbf{antecedente dominial de la matrícula 1.539, Finca San Alejo}, y desde cuándo pertenece a los titulares que constituyeron \textsc{agro san alejo s.a.s.} en enero de 2026 \textit{(cap.\ \ref{cap:trabajo})}
\item \textbf{padrón catastral correspondiente al Revalúo General de 1944}, que permitiría convertir las cuarenta y siete valuaciones publicadas en superficies y en ubicaciones, y establecer si esas partidas son todas las del departamento o sólo las revaluadas. \textbf{Interesa además el domicilio declarado de cada titular}: es el único dato que permitiría medir qué proporción de la propiedad del departamento estaba, en 1944, en manos de residentes de la capital, y de él depende la hipótesis sobre el veraneo que el capítulo \ref{cap:poblacion} formula y no resuelve \textit{(caps.\ \ref{cap:fincas} y \ref{cap:poblacion})}
\item expediente y \textbf{plano de mensura de la finca ``La Helvecia'', antes ``Calderilla''} (Juzgado de 1ª Nominación en lo Civil, auto del 10/04/1937, perito Jorge de Bancarel): tiene por límite occidental el río La Caldera y por lindero norte a quien veinte días después presidiría el municipio \textit{(cap.\ \ref{cap:fincas})}
\item \textbf{folio 374, asiento 475 del Libro D de Títulos de La Caldera}: la inscripción de la finca \textbf{Campo Alegre} a nombre de Isaac Fernández, por compra a Teodoro Sosaya en noviembre de 1937. Permite reconstruir la cadena de titularidad hasta la expropiación de 1972 y establecer quién era dueño de la tierra del embalse cuando el Estado decidió tomarla \textit{(cap.\ \ref{cap:expropiacion})}
\item \textit{(a pedir también al municipio y al Archivo General)} \textbf{Ordenanza del Municipio de La Caldera del 13 de octubre de 1914}, que reglamenta la distribución de la totalidad del caudal del río Santa Rufina y reconoce derechos adquiridos por uso tradicional: es el instrumento de gestión hídrica más antiguo del departamento que el archivo menciona. No está en ninguna de las setenta ediciones de 1914 ni en las cuarenta y nueve de 1915 que este relevamiento leyó una por una. \textbf{La vía más corta no es el municipio: son las fojas 19-20 y 69-72 del expediente 8312-P/937}, donde el decreto de 1939 sitúa el disenso municipal que invocó la ordenanza ---para invocarla, el municipio tuvo que acompañarla o transcribirla---. En el municipio, lo que corresponde pedir de esa época no es un boletín, que no existía, sino el \textbf{libro de actas de la comisión municipal} \textit{(cap.\ \ref{cap:expropiacion})}
\item expediente \textbf{8312-P/937} completo, con el disenso municipal de fs.\ 19-20 y 69-72 y las oposiciones particulares de fs.\ 26-28, 38-41 y 50-51 \textit{(cap.\ \ref{cap:expropiacion})}
\item expediente \textbf{8312-F} ---que el decreto de 1939 cita como 8312-P/937---: si es el mismo trámite, y qué pasó con los usos tradicionales del río Santa Rufina al construirse el embalse \textit{(cap.\ \ref{cap:expropiacion})}
\item \textbf{informe y antecedentes de la Comisión ad-hoc del Decreto 14091 del 23 de octubre de 1931} (expte.\ 4920-M), sobre las concesiones de agua de los ríos Vaqueros y Wierna y las ordenanzas municipales que las reglaban: es el único estudio integral de las concesiones de esta cuenca que el Estado ordenó hacer, y este libro no encontró rastro de su resultado \textit{(cap.\ \ref{cap:redes})}
\item \textbf{Ordenanza General de Impuestos de la Municipalidad de La Caldera vigente en 1930}, que la Resolución 434 del Ministerio de Gobierno cita y verifica sin transcribir \textit{(cap.\ \ref{cap:hacienda})}
\item \textbf{expedientes 1766-M/941, 190-M/943, 1549-M/943, 6463-M/944 y 8266/944}, los cinco devueltos por sus propios decretos a la Municipalidad de La Caldera: contienen la Ordenanza General de Impuestos y los presupuestos municipales de 1943, 1944 y 1945 y los balances de 1940 y 1942, ninguno de los cuales se publicó. El presupuesto de 1945 está, según el art.\ 1º del Decreto 5786-G, \textbf{a fojas 19 y 20} del último \textit{(cap.\ \ref{cap:hacienda})}
\item \textbf{plano adjunto al expediente 17.870/1944} de la Dirección General de Hidráulica, al que remite el art.\ 2º del Decreto 4614-H: es lo único que permite ubicar la toma de los Ferrocarriles del Estado sobre el río Mojotoro \textit{(cap.\ \ref{cap:aguabaja})}
\item \textbf{folio 342, asiento 438 del Libro B de La Caldera}: la inscripción de la finca Chalchanio, citada por el edicto de 1937 \textbf{y de nuevo por el remate de 1945 sobre una cuarta parte indivisa}, con otro ejecutado ---de modo que el asiento registra al menos dos traslaciones y una partición entre comuneros, uno de ellos Ernesto Echenique. Es, con el folio 374 del Libro D de Campo Alegre, una de las dos referencias registrales concretas que el archivo entrega para fincas del departamento, y permitiría verificar si el ejecutado Ramón I.\ Juárez es el comisario de policía homónimo \textit{(cap.\ \ref{cap:fincas})}
\item antecedente dominial de la finca \textbf{``Vaqueros o Entre Ríos''} y resultado del remate de mayo de 1935 en la ejecución del Banco Provincial contra Ramón Bracheri: es la finca que da nombre al municipio \textit{(cap.\ \ref{cap:fincas})}
\item \textbf{planillas y planos de la catastración de 1929} del departamento La Caldera, levantada por Adeodato Aybar\index{Aybar, Adeodato} (expediente 3463-D de la Comisión de Catastro): son el registro más antiguo de la propiedad departamental y el punto de partida natural de cualquier cadena de títulos \textit{(cap.\ \ref{cap:opacidad})}
\item \textbf{matrícula y antecedente dominial de dos fraccionamientos del paraje El Gallinato}, identificados sobre el visor de IDESA y sin datos asociados \textit{(cap.\ \ref{cap:opacidad})}: \textbf{(a)} el conjunto de parcelas en cintas perpendiculares al arroyo, al este del cerro Pucheta, en torno a \textbf{--24,6802 / --65,3103}; y \textbf{(b)} el conjunto de parcelas paralelas al valle, al norte del anterior. Interesa \textbf{cuándo se subdividió cada uno, por qué instrumento y a partir de qué matrícula de origen}
\item antecedentes dominiales de las fincas El Gallinato, San Alejo, Los Porongos, San
Ignacio, Las Cañas y Potrero de Castilla; correspondencia entre los trece partidos de 1909
y las matrículas actuales, en la Dirección General de Inmuebles; y fecha de la donación
Serrey\index{Serrey, Carlos}, fecha de la sentencia sucesoria, superficie del dominio partido y cadena de
titularidades resultante \textit{(cap.\ \ref{cap:fincas})}
\item estado registral de la matrícula 1.303: si el IPDUV construyó viviendas, si hubo reversión al dominio provincial, y qué hay hoy en ese terreno \textit{(cap.\ \ref{cap:expropiacion})}
\item \textbf{antecedente dominial de las matrículas 4366 y 6.558}, para establecer si la segunda proviene de la primera; \textbf{expediente 34-167.127/13 completo}, con el peticionante, que contiene las tres resoluciones y los informes técnicos de la comisión; \textbf{expediente 34-253.934/13}, de las defensas con gaviones a las que la Resolución 023/14 sujetó la línea de ribera del pueblo, y si se construyeron \textit{(el anexo de la 023/14 ya se obtuvo)}; aclaración sobre el expediente 34-167.127/13 y a qué resolución corresponde; verificación de las coordenadas de las Resoluciones 6/14 y 7/14; y las determinaciones resultantes de las comisiones aprobadas por las Resoluciones 395/13 y 8/14 \textit{(cap.\ \ref{cap:ribera})}
\item fojas 58, 59, 68, 84, 97/101 y 137 del expediente 18-37.348/2019-0 y el expediente 18-37349/19-1; plano 1.164 completo, con sus tres láminas ---la Resolución Delegada 646/21 ya se obtuvo---; situación de las Etapas 3 y 5; ordenanza o acto municipal de aprobación anterior a 2019, si existe; y legajo societario y composición actual de Jardín Celestial S.R.L. en la Dirección General de Sociedades ---el antecedente dominial de 4.366 y 6.558 figura en el ítem anterior--- \textit{(cap.\ \ref{cap:loteo})}
\item expediente 18-27.372/11 ---con el plano de mensura y subdivisión, para explicar los 8.622,33 m² que el balance publicado no distribuye entre sus rubros--- y el Reglamento de Condominio de Las Vertientes ---la resolución de línea de ribera de ese loteo, la 80/14, ya se obtuvo---; texto de la Ley 5602 y del Decreto 924/80 sobre clubes de campo; antecedentes dominiales de la matrícula 3989 fracción 78; y estado actual del loteo: lotes vendidos, edificados y su fuente de agua \textit{(cap.\ \ref{cap:loteo})}
\item expediente 0110018-20984/07-0 de Chacras del Gallinato y estado actual del proyecto; \textbf{legajo societario de El Gallinato S.R.L.}; decreto de concesión de agua pública subterránea exigido por el art.\ 3º inc.\ f; expedientes del Barrio La Misión y \textbf{plano con el retiro por línea de ribera}, que documentaría cómo se grafica ese descuento; y antecedentes dominiales de las matrículas 3.709, 4.065 y su relación con la Finca El Gallinato \textit{(cap.\ \ref{cap:loteo})}
\item \textbf{folio real de las matrículas 4097 y 6663}: la aparición de la 6663 junto a la 4097 en la determinación de 2022 sugiere una subdivisión que sólo el folio acredita \textit{(cap.\ \ref{cap:ribera})}
\item expediente 01201.59-198642/2012-0 --- el \textbf{plano de mensura Nº 000962} ya se obtuvo y se reproduce en la lámina \ref{lam:plano962} ---; antecedentes dominiales de las matrículas \textbf{3.160, 3.161 y 3.162} y su relación con la Finca El Gallinato del artículo 6º del PDUA ---el de la 3.709 figura en el ítem de Chacras del Gallinato---; y nómina de prescripciones adquisitivas administrativas declaradas a favor de la Provincia en el departamento \textit{(cap.\ \ref{cap:tierrafiscal})}
\item \textbf{informes de la Unidad de Seguimiento y Control del Servicio de Información Catastral} previsto en las actas del Decreto 1639/07 con La Caldera y Vaqueros: mejoras detectadas en cada municipio, montos percibidos por la Provincia y prórrogas del convenio; y \textbf{destino del proyecto de ley de 2007} que debía regular el servicio para todos los municipios \textit{(cap.\ \ref{cap:hacienda})}
\end{itemize}

\section{Secretaría de Recursos Hídricos}

\begin{itemize}[leftmargin=*]
\item \textbf{expedientes 34-21.387/12 y 34-129.165/10} completos ---el primero contiene el aviso de inicio de abril de 2012 y la rectificación de marzo de 2013; el segundo, la determinación con coordenadas de diciembre de 2012---, para establecer \textbf{por qué una misma determinación sobre el arroyo Guaranguay tramitó en dos expedientes distintos} \textit{(cap.\ \ref{cap:ribera})}
\item \textbf{expedientes de los consorcios de agua pública} 0090034-86490 ---Loteo Valle Alegre--- y 0090034-181690/2020-0 ---Valle Hermoso / Coronado Cueto---: acto de constitución, estatuto, área demarcada, concesiones que los amparan y padrón de usuarios \textit{(cap.\ \ref{cap:aguabaja})}
\item \textbf{expediente 34-2.253/56}, de la finca «Vaqueros», iniciado en 1956 y aún publicándose en 2010: \textbf{qué se resolvió y cuándo} \textit{(cap.\ \ref{cap:aguabaja})}
\item \textbf{expediente 0050034-15.196/10 de El Gallinato S.R.L.}, de concesión de agua pública: caudal, uso, fuente y resolución que lo cerró \textit{(cap.\ \ref{cap:aguabaja})}
\item \textbf{si el geólogo «Víctor Omar Viera» del edicto de octubre de 2011 y el «Víctor Omar Nieva» de los de 2012 son la misma persona} \textit{(cap.\ \ref{cap:ribera})}
\item \textbf{determinaciones de línea de ribera del departamento anteriores a 2004}, si las hay: los avisos más antiguos que la serie registra son de 2004 y 2005, y el Decreto 1.989/02 rige desde 2002 \textit{(cap.\ \ref{cap:ribera})}
\item \textbf{texto de la Resolución Nº 417 del 02/09/2011} y su expediente 34-50.034--4.034/11, para establecer \textbf{si el acto dice «río Arenales» como su edicto} o si el error es sólo de la publicación; y \textbf{si se rectificó} \textit{(cap.\ \ref{cap:ribera})}
\item \textbf{nómina de pedidos de rectificación de línea de ribera} sobre los ríos del departamento, con su resultado, e \textbf{inventario de Certificados de No Inundabilidad} emitidos para loteos de La Caldera y Vaqueros, con sus constancias de pago de arancel: la propia Secretaría describe la rectificación como práctica habitual y conflictiva, y el certificado como trámite necesario, y en los expedientes relevados sólo consta un certificado de no inundabilidad, el que presentó Ayres de Quitilipi en 2020 \textit{(cap.\ \ref{cap:ribera})}
\item \textbf{series de las dos estaciones activas del río Mojotoro} --- «Mojotoro -- Güemes» de la Red Hidrológica Nacional (\textit{sitecode} 2103, base BDHI) y su homónima automática de la red de alerta (\textit{sitecode} 6365) ---: desde cuándo operan, con qué continuidad, y si alguna determinación de línea de ribera del departamento las invocó \textit{(cap.\ \ref{cap:ribera})}
\item \textbf{resultados del Plan de Monitoreo de calidad del agua de la subcuenca Mojotoro} desde 2012 ---siete puntos en la microcuenca del río La Caldera---, y el sentido de la «Línea Ribera --- Res.\ 236--360 SRH» que rotula el plano de Aguas del Norte de 2022: si la 236 es la 236/14, que rectificó la 142/14 \textit{(cap.\ \ref{cap:amparo})}; y \textbf{series de caudales de las cinco estaciones de aforo} --- El Volcán (Las Nieves), El Angosto (Mojotoro), San Alejo, Santa Rufina y Desembocadura al Nieves (Yacones) ---, publicadas en la \textbf{Estadística Hidrológica de la República Argentina (2004)} y en el \textbf{Atlas Digital de los Recursos Hídricos Superficiales} del Instituto Nacional del Agua. \textbf{No es un pedido de archivo: son publicaciones editadas}, citadas al pie de las tablas de estaciones del Plan de Manejo. Las cinco series de caudal medio mensual ya se obtuvieron del SNIH; faltan los caudales máximos instantáneos de las cinco y la explicación de los cuarenta y cuatro meses faltantes de El Angosto, concentrados en 1967--1969 y 1975--1982 \textit{(cap.\ \ref{cap:ribera})}
\item \textbf{serie completa del Observatorio Pluviométrico Los Yacones}, en funcionamiento con observador pago al menos en 1980, y nómina de estaciones de la red provincial en la cuenca del río La Caldera: es el único dato instrumental de origen provincial que este relevamiento encontró \textit{(cap.\ \ref{cap:defensas})}
\item \textbf{expedientes 181-C, 386-R y 2225-O de 1935 y 68-D de 1938}, con los \textbf{croquis y cómputos métricos de cada obra}, que permitirían establecer si las defensas de la margen derecha desplazaron el problema a la izquierda; el de diciembre de 1935, además, con el \textbf{levantamiento de la zona, los planos y el presupuesto} del proyecto de defensas elevado por Obras Públicas en diciembre, con el croquis del espigón de emergencia y la rendición de cuentas de Augusto Regis\index{Regis, Augusto} ante Contaduría General; y \textbf{cualquier estudio posterior sobre el ``espigón macizo''} que ese decreto consideró la solución definitiva y que debía estudiarse, aunque no lo encargó \textit{(cap.\ \ref{cap:defensas})}
\item \textbf{qué pasó después del fallo de 1931}: si la concesión de cien litros por segundo del río Wierna se restableció efectivamente, ---la acción de daños se promovió y la Corte la rechazó en 1936---, y si esa concesión subsiste hoy sobre el inmueble que en 2014 llegó a audiencia pública como Urbanización Cerro de Buena Vista \textit{(cap.\ \ref{cap:politica})}
\item documento fuente de la serie 229 m (2019) --- 89 m (2021) del cauce del río Vaqueros, y su metodología; expedientes de corrimiento de línea de ribera tramitados en el departamento; permisos de extracción de áridos vigentes sobre el río Vaqueros y el río La Caldera; y superficie efectivamente cedida como espacio verde en los loteos aprobados, contrastada con su ubicación respecto de la línea de ribera \textit{(cap.\ \ref{cap:vaqueros})}
\item expediente 7766/15 y nómina completa de intervenciones mecánicas sobre el cauce del río La Caldera entre 2013 y 2026, por organismo, con sus autorizaciones ambientales e hídricas si las hubo; y norma que rija la evaluación previa de las obras del propio Estado sobre cauces \textit{(cap.\ \ref{cap:defensas})}
\item evaluaciones del art.\ 171 previas a los permisos de extracción de áridos sobre el río La Caldera y sus arroyos; y actuaciones del art.\ 173, si las hubo \textit{(cap.\ \ref{cap:ribera})}
\item \textbf{Certificado de No Inundabilidad} del loteo El Durazno, o constancia de su inexistencia; \textbf{carátula del plano 1.164}, para verificar si menciona el número de la resolución de línea de ribera; expediente 0090034-242142/2024-0 y Resoluciones 279/17, 306/19 y 108/2025 sobre el río La Caldera; e integración nominal de las Comisiones Técnicas intervinientes \textit{(cap.\ \ref{cap:ribera})}
\item organigrama provincial y ley de ministerios: el cambio de rango \textbf{ya está acreditado} ---en julio de 2025 los edictos los firma la Secretaría de Recursos Hídricos, y el 13 de febrero de 2026 el Decreto 52/26 ya designa a Romero Leal Subsecretario \textit{(cap.\ \ref{cap:ribera})}---, y la Decisión Administrativa 569/25 lo ubica dentro de la Secretaría de Obras Públicas; lo que resta pedir es \textbf{el manual de misiones y funciones de la Subsecretaría y qué competencias conserva}, dado que la Ley 8.511 asigna los recursos hídricos al Ministerio de Producción y Minería; y presupuesto asignado a la Subsecretaría en los ejercicios 2019 a 2026, para contrastar con su alegación de falta de capacidad \textit{(cap.\ \ref{cap:amparo})}
\item \textbf{expediente 0090034-211560/2014-0} completo, con las características técnicas del dren que el propio aviso pone a disposición, y la resolución que otorgó o denegó la concesión; \textbf{antecedentes dominiales del catastro 4366}, con la nómina de co-titulares registrales y sus porcentajes; y régimen aplicable a los drenes horizontales de subálveo, para establecer si el art.\ 143 los alcanza \textit{(cap.\ \ref{cap:loteo})}
\item coordenadas Posgar de la \textbf{Resolución 360/16} (río Wierna, expte.\ 34-171.172/15) y relación de sus siete matrículas con las canteras del Wierna y con las obras de encauzamiento sobre ese río \textit{(cap.\ \ref{cap:ribera})}
\item \textbf{resolución, instructivo o criterio escrito que fije cuándo un edicto de línea de ribera publica las coordenadas y cuándo las remite a la sede de la Secretaría}; y \textbf{coordenadas de las Resoluciones 142/14 y 236/14} sobre la matrícula 3616, ni las erróneas ni las corregidas publicadas; y \textbf{qué línea de ribera rige para las matrículas 3593 y 3584 de las Resoluciones 6/14 y 7/14}, cuyos avisos publicaron la tabla de la línea del pueblo que después aprobó la 023/14 \textit{(caps.\ \ref{cap:planteo}, \ref{cap:ribera} y \ref{cap:opacidad})}
\item texto íntegro de la Resolución 383/14 y las coordenadas Posgar de su línea de ribera, para establecer si determinó zona inundable y si invocó el art.\ 172; autorizaciones del art.\ 172 para construcciones en zona ribereña de los arroyos El Durazno y Guaranguay; superficie del loteo comprendida en zona inundable; y relación dominial entre las matrículas 1590, 4146, 4366 y 6.558 del departamento \textit{(cap.\ \ref{cap:ribera})}
\item \textbf{edicto inicial de la determinación de línea de ribera de los arroyos El Durazno y Guaranguay}, presumiblemente publicado entre 2013 y 2014, con \textbf{la resolución que aprobó la Comisión Técnica para esos dos arroyos} ---o constancia escrita de que no existe: dentro del expediente 34-167.127/13, la Resolución 395/13 sólo nombra al río La Caldera y al arroyo El Manzano, y el buscador oficial de avisos no devuelve ningún acto que nombre a los otros dos; interesa \textbf{cuándo y por qué se incorporaron al trámite}---, sus integrantes, el método elegido y la prohibición cautelar; asiento en el Libro de Catastro de Aguas; y constancia de la notificación a la Dirección General de Inmuebles y de su uso en el visado del plano 1.164 \textit{(cap.\ \ref{cap:ribera})}
\item \textbf{expediente 0090034-25.217/2014-0/3}, con las características técnicas de la captación que el edicto pone a disposición; expediente 0090034-168.839/14 y antecedentes del catastro 2907 y del catastro de origen 126, para establecer qué subdivisión originó el traslado de la asignación de riego; titularidad y ubicación del Catastro 1700 de La Caldera, y si corresponde a un loteo; resolución que otorgó o denegó la concesión; y nómina completa de concesiones de agua pública otorgadas en el departamento \textit{(cap.\ \ref{cap:redes})}; y \textbf{cuadro completo de observaciones al padrón del Consorcio de Usuarios «Mojotoro y Regantes del Río La Caldera»}, elevado por la Dirección General de Concesiones Hídricas el 2 de octubre de 2023 \textit{(cap.\ \ref{cap:amparo})}
\item expediente 34-6147/05 y resolución final de la concesión de Finca Lesser; \textbf{juicio sucesorio de Adolfo Saravia Valdez}: juzgado, número de expediente y autorización judicial para desarrollar el loteo; antecedentes dominiales del catastro 1869; caudal exacto autorizado para uso recreativo, sobre el original impreso; y nómina de desarrollos bajo la figura de club de campo en el departamento \textit{(cap.\ \ref{cap:redes})}
\item \textbf{Resolución 000056/2025} y expediente 0090034-185051/2024-0, con los puntos
de vinculación de la línea de ribera del arroyo Vaqueros a la latitud del catastro 1122; y
serie completa de determinaciones sobre el departamento \textbf{entre 2016 y 2026}, que este
relevamiento no cubre \textit{(cap.\ \ref{cap:ribera})}
\item \textbf{la resolución de determinación de línea de ribera del río La Caldera en la
zona de emplazamiento de la planta de tratamiento del Acueducto Norte}, posterior a agosto de
2017 ---podría ser la Resolución 279/17, sobre el río La Caldera, matrícula 4146, cuya comisión se designó por la 291/16---, y los \textbf{requisitos fijados para los cruces de los cauces de Los Naranjos,
Guaranguay y Wierna}; y el expediente de la obra del Acueducto Norte en lo relativo a esas
condiciones hídricas \textit{(cap.\ \ref{cap:ribera})}
\item las \textbf{capas vectoriales} de los tres mapas temáticos de mayo de 2022 --- las
ocho hojas del informe ya se obtuvieron y se reproducen en las láminas \ref{lam:mapa1} a
\ref{lam:mapa3} y \ref{lam:mojotoro} ---; los polígonos faltantes del cuadro de subcuencas
(identificadores 3, 6 y 10 a 13); y los dos estudios inéditos de 2008 y 2009 que el informe
declara entre sus fuentes \textit{(cap.\ \ref{cap:amparo})}
\item \textbf{anexo de montos por municipio} del Plan de Mínimas y Encauzamiento
2016--2017, y sus \textbf{certificados de avance} para La Caldera y Vaqueros: el Acta Acuerdo
fija el mecanismo y no las cifras, y el Decreto 1.889/16 publica el acta pero no el detalle;
\textbf{el expediente de gestión de fondos iniciado el 31 de julio de 2020} y su resultado; y
el \textbf{instrumento que ampara la colaboración de Supercemento} con maquinarias en las obras
municipales de 2020, si existe \textit{(cap.\ \ref{cap:amparo})}
\item \textbf{solicitantes de los cuatro expedientes} del arroyo Chaile en la manzana 148 de Vaqueros \textit{(cap.\ \ref{cap:ribera})}
\item \textbf{expediente 0090034-160869/2013}, con el informe de la comisión de 2013 y los fundamentos de la Resolución 142/2025 que la cierra \textit{(cap.\ \ref{cap:ribera})}
\item \textbf{fecha de la Resolución 111/2025} en el registro de la Secretaría, que el edicto atribuye al mismo día que la 142/2025 \textit{(cap.\ \ref{cap:ribera})}
\item \textbf{textos de los avisos de línea de ribera no leídos}: expedientes 34-3.998/03 (2004), catastro 2961 y Resolución 683 (2005), fundamentos de la ampliación de la Resolución 414/16 (río Vaqueros) por la 144/17, y \textbf{fecha de la 414/16}, que la 144/17 da como 5 de noviembre de 2016 y la numeración ubica después del 21 de noviembre; aclaración de si el expediente 34-125.733/15 recae sobre el río Trancas o el Wierna; determinaciones resultantes de las Resoluciones 122/16 (arroyo Quintín, matrícula 4953) y 290/16 (arroyo Chaile, matrícula 1534), y relación entre la comisión de la 122/16 y la determinación «por Administración» de la \textbf{Resolución 389/16}; expediente correcto de la \textbf{Resolución 131/16}; comisión de la \textbf{Resolución 282/16}; solicitante de la Resolución 197/15 (catastro 4049, río La Caldera, aviso titulado con Aguas del Norte) y acto que aprobó esa determinación; zona inundable de la \textbf{Resolución 44/2021} (matrícula 4097); aclaración del número y del curso del expediente de la matrícula 2020; textos de la \textbf{Resolución 009/12} y de su rectificación por la \textbf{130/2025} (matrícula 216); relación entre las dos determinaciones de la matrícula 3017 (Resoluciones 125/2021 y 107/2022); revisión de las páginas restantes del listado de avisos de 2015 a 2017; y búsqueda de los avisos posteriores a 2017 con otros términos \textit{(cap.\ \ref{cap:ribera})}
\item \textbf{identificación del inmueble «s/Nº» de la sección B de Vaqueros} sobre el que la Resolución 201/13 determinó la línea de ribera del río Vaqueros, y su relación con la matrícula 40 de las Resoluciones 414/16 y 144/17, cuya línea continúa la de 2013 \textit{(cap.\ \ref{cap:ribera})}
\item \textbf{resolución que fija la dotación de riego de 0,525 litros por segundo por hectárea} que todos los edictos de subdivisión aplican y que ya aplican los reconocimientos de 1948 \textit{(cap.\ \ref{cap:siglo})}: número, fecha, base de cálculo y fuentes a las que se aplica \textit{(cap.\ \ref{cap:aguabaja})}
\end{itemize}

\section{Secretaría de Tierras y Bienes del Estado}

\begin{itemize}[leftmargin=*]
\item texto de la \textbf{Resolución 3/16} de la Secretaría, comprendida en la misma fe de erratas que las Resoluciones 7, 8 y 9 de 2016 (B.O.\ Nº 19.911) y ausente de este relevamiento \textit{(cap.\ \ref{cap:tierrafiscal})}
\item \textbf{ubicación y forma de acceso al registro público de Tenedores Precarios y Adjudicatarios} que ordena llevar el art.\ 1º inc.\ g; resoluciones dictadas en consecuencia del art.\ 1º inc.\ f; sorteos efectivamente realizados desde 2016 en el departamento La Caldera y en Vaqueros, con sus actas notariales; y convenios celebrados con ambos municipios conforme el art.\ 1º inc.\ c \textit{(cap.\ \ref{cap:tierrafiscal})}
\item expediente 388-214146/24 y el informe del área competente, en particular sobre la conversión a «sin cargo» del segundo caso del Barrio Martel de Tartagal (Decretos 2067/05 y 3354/12) y sobre la discordancia entre las Resoluciones 156D/23 y 159D/23; expediente 131-39395/04 del Barrio Leopoldo Lugones, para establecer si el error de parcela y matrícula tuvo consecuencia práctica; y art.\ 76 de la Ley 5.348/78 y Decreto 41/96 con sus modificatorios \textit{(cap.\ \ref{cap:tierrafiscal})}
\item expedientes 0010226-19595/2010-0 y concordantes; \textbf{plano de mensura y desmembramiento Nº 000943} de las cuarenta y una hectáreas; \textbf{legajo societario de Proyecto 1 S.A.}; estado de cumplimiento del plazo de cuatro años para construir, vencido en 2019, y si la Provincia verificó o dispuso la restitución; antecedente dominial de la matrícula 4.425 y motivo por el cual estaba en el dominio privado del Estado; y nómina de comodatos de tierra fiscal otorgados en el departamento \textit{(cap.\ \ref{cap:tierrafiscal})}
\end{itemize}

\section{Juzgado de Minas y de Registro}

\begin{itemize}[leftmargin=*]
\item \textbf{expedientes de la mina de plomo «20 de Febrero»}, ubicada en 1918 en la finca «Cerro Nevado» de Garzón y Pintos, y de la mina \textbf{«Sara»}, denunciada por despueble el mismo año, las dos por Miguel A.\ Sosa y Dermidio Arancibia; y \textbf{ubicación de la finca Cerro Nevado} en el catastro actual \textit{(cap.\ \ref{cap:siglo})}

\item \textbf{expediente de la cantera «Macarena»} de \textsc{transporte hnos.\ mendoza s.r.l.}, río Wierna, diez hectáreas fiscales: \textbf{volumen autorizado de extracción, canon abonado e informe de impacto ambiental}, y \textbf{a quién se vende el material} \textit{(cap.\ \ref{cap:infraestructura})}
\item \textbf{si \textsc{cantera el rescoldo s.a.} tramitó concesión de áridos} ante el Juzgado de Minas, y en su caso sobre qué curso y con qué informe de impacto ambiental \textit{(cap.\ \ref{cap:trabajo})}
\item constancia de que el \textbf{remate administrativo convocado para el 17 de julio de 1941} (Decreto 4769, B.O.\ Nº 1.899) de las minas «Caldera Nº 1, 2 y 3» quedó desierto ---lo indica la anulación de patentes de 1943---, y si la finca San José conserva pertenencias mineras vigentes \textit{(cap.\ \ref{cap:siglo})}
\item expedientes 1168-C ---cateo de Larrán de 1925--- y 44-L: estado de las minas de plomo y plata \textbf{Caldera Nº 1, 2 y 3} sobre la finca San José, y cadena de titulares de esa finca desde los Quipildor Hermanos; y \textbf{expediente 38-H} (cateo de Alberto Jaime Harrison y Juan Larrán sobre Las Nieves, pedido en 1929 y concedido en 1930), para establecer qué relación tuvieron con la manifestación 44-L, registrada el 30 de abril de 1930 \textit{(cap.\ \ref{cap:siglo})}
\item límite de superficie \textbf{por solicitante} en el Código de Minería vigente en 1925, y expedientes 1163 C, 1164 C, 1165 C y 1167 C, presentados en diez días por cuatro personas de apellido Prémoli con socios distintos \textit{(cap.\ \ref{cap:siglo})}
\item expediente 22.295 del Juzgado de Minas: estado actual de la mina Cristina M-1, oposiciones si las hubo, y evaluación ambiental si correspondía; \textbf{legajo societario de TURPINAR S.A.}\ y superficie total de las matrículas 1.981 a 1.985; titularidad de la matrícula 1.862, sujeta a inscripción registral; y nómina completa de manifestaciones de descubrimiento y minas en el departamento \textit{(cap.\ \ref{cap:resistencias})}
\item expedientes 18.219 y 19.199/08 del Juzgado de Minas: evaluación del art.\ 171 del Código de Aguas, canon, plazo de las concesiones y volumen autorizado; nómina completa de canteras de áridos concesionadas en el departamento y su relación con la línea de ribera; e identificación catastral de las hectáreas fiscales afectadas por ambas \textit{(cap.\ \ref{cap:resistencias})}
\item texto íntegro de la sentencia de 2021; texto de la del 7 de julio de 2023; \textbf{anexos del informe final del taller participativo} (el informe ya se obtuvo), diseñado para el 29 de septiembre de 2023, la desgrabación y la constancia de la devolución pública; actas de asamblea de las comunidades Kolla Kondorwaira y Kolla Urkuwasi sobre la consulta de septiembre de 2023, si las hubo: el dictamen provincial del 4 de septiembre las recomendó y la actuación 9.947.016, con la que el juzgado dio por acreditada la consulta, no las incluye ---se pide a las propias comunidades---; las presentaciones escritas de las comunidades kolla sobre obras hechas sin autorización ni consulta; y el padrón completo de regantes del departamento que acompañó el Dictamen 129/2023 \textit{(cap.\ \ref{cap:amparo})}; \textbf{ejemplar del Plan de Manejo Integral aprobado que quedó a disposición de la comunidad en Avenida General Güemes s/n, medidor 994, La Caldera}, para cotejarlo con el acompañado en autos que este libro leyó; ordenanza que establece el Protocolo de Crecidas e Inundaciones y crea el Comité Operativo de Emergencia \textit{(cap.\ \ref{cap:amparo})}
\item expediente completo y su incidente INC 23860/1: acta de citación de Jardín Celestial S.R.L., su carácter procesal y sus presentaciones; providencias del Expte.\ 780280 de ejecución; requerimiento a la Municipalidad de Vaqueros y consecuencias de su omisión; actuación de la Fiscalía Civil, Comercial y Laboral Nº 2 desde marzo de 2024 \textit{(cap.\ \ref{cap:amparo})}
\item \textbf{expediente MIN 780280/22} (17-00780280-0), Juzgado de Minas, Distrito Centro. En orden de prioridad: la \textbf{resolución interlocutoria del 12 de febrero de 2026} en su texto de actuación (16.465.143; el acta de la audiencia, con lo resuelto, ya se obtuvo); los \textbf{anexos cartográficos del Plan de Manejo} en su formato original (actuaciones 13.175.771, 13.175.751 y 13.196.832); la \textbf{Línea de Base Socio Ambiental} de noviembre de 2022; los \textbf{informes de los veedores}; y las \textbf{presentaciones posteriores al 12 de febrero de 2026} ---planificación de obras, pedido y concesión de prórroga, informes de cumplimiento---. Todo el expediente es consultable por el sistema público de expediente digital \textit{(cap.\ \ref{cap:amparo})}
\item \textbf{legajos societarios de Dal Borgo Construcciones S.R.L.\ y de Ingeniero Medina S.A.}: socios y administradores, para establecer si César Domingo Dal Borgo y Héctor Enrique Medina ---concesionarios de las canteras «La Vaquera», «La Generosa» y «Tomasito» y de una servidumbre de acopio--- integran las sociedades que ofertaron en la licitación del puente sobre el río Vaqueros \textit{(cap.\ \ref{cap:redes})}
\item \textbf{formularios de las encuestas} del diagnóstico social del Plan de Manejo, realizadas entre febrero y julio de 2024 y conservadas por YSATI Ingeniería S.R.L.\ «a disposición del juzgado»: son el material de campo que este libro no tiene, y su consulta exigiría resguardar el anonimato bajo el cual se tomaron \textit{(cap.\ \ref{cap:amparo})}
\item las \textbf{capas vectoriales del catastro minero} de la cuenca y el
\textbf{estado actualizado} de los informes de impacto ambiental de las veinte concesiones ---
los tres cuadros de 2021, 2020 y 2018 ya se obtuvieron ---, en particular si las seis que
constaban como no aprobadas siguen vigentes; y las \textbf{actas de las inspecciones
periódicas} que el Programa declara realizar sobre las canteras del departamento \textit{(cap.\ \ref{cap:resistencias})}
\end{itemize}

\section{Municipios de La Caldera y Vaqueros}

\begin{itemize}[leftmargin=*]
\item \textbf{cálculo de recursos de la Comisión Municipal de La Caldera para 1917} y actuaciones de los decretos 1.557 y 1.660 de 1918: es la planilla en la que \$3.083,94 de saldo del ejercicio anterior se computaron como renta del año, y por cuya depuración el departamento perdió la municipalidad electiva siete semanas después de obtenerla. \textbf{La vía más corta es el libro de actas de la comisión municipal}, que este apéndice ya reclama para 1914 \textit{(cap.\ \ref{cap:siglo})}
\item \textbf{notas del Comisionado Municipal de La Caldera de septiembre y noviembre de 1918} proponiendo los jueces de agua de los distritos de Caldera y Calderilla, que los decretos 113 y 190 citan como su antecedente y que no se publican: son el último registro conocido de la facultad municipal sobre el agua antes de 1942 \textit{(caps.\ \ref{cap:planteo} y \ref{cap:redes})}

\item \textbf{expedientes municipales sin respuesta}: informe de impacto ambiental para intervenir el río en el área de exclusión (90302-250447/2023), nota a Minería del 10/07/2025 (302/138756), proyecto de defensas de Santiago Apóstol y Nogalar 4 (386-219264), y pedidos de fondos para honorarios (011366-142874/2025); y estado de la \textbf{ejecución de honorarios 938.744/25} (\$10.004.770, costas solidarias de Provincia y municipio) y del embargo sobre la matrícula 1.846, donada por la Ley 7699; y \textbf{ordenanza que aprobó el protocolo de emergencia reformado} en marzo de 2024 \textit{(cap.\ \ref{cap:amparo})}
\item \textbf{acta de la multa ambiental} aplicada al emprendimiento que habría desmontado monte ribereño sin autorización, \textbf{expediente sancionatorio} de la Secretaría de Ambiente, y \textbf{versión taquigráfica de la audiencia pública} municipal en la que ese proyecto se trató: son los tres documentos que resolverían las denuncias del Centro Vecinal de La Calderilla de enero de 2026 \textit{(cap.\ \ref{cap:infraestructura})}
\item \textbf{balances de fondos de la Comisión Municipal de La Caldera} de los ejercicios 1908 a 1940, que el Poder Ejecutivo aprobaba por decreto: permitirían reconstruir la serie completa de la renta municipal, de la que este libro sólo tiene dos puntos \textit{(cap.\ \ref{cap:hacienda})}
\item resultado de las licitaciones de 1935 para la adquisición de terrenos destinados a edificación escolar en Vaqueros y a asistencia social en La Caldera: qué inmuebles se compraron, a quién y en cuánto \textit{(cap.\ \ref{cap:tierrafiscal})}
\item identificar el emprendimiento referido, el expediente sancionatorio y el acta de la sesión del Concejo; y, como el buscador de avisos del Boletín Oficial no devuelve convocatoria a audiencia pública en el municipio de La Caldera entre 2019 y 2026, \textbf{el decreto u ordenanza municipal que la convocó, si existió}, con su fecha, y \textbf{quién presidía el Concejo en esa fecha} (actas de sesiones preparatorias): la prensa oficial registra como presidentas a Yamila Sastre (2021) y a Ana Sumbay (sin fecha), y como presidentes a Germán Sumbay (enero de 2025 y enero de 2026) y a Jorge García (agosto de 2026), de modo que la audiencia sería anterior a 2025 \textit{(cap.\ \ref{cap:infraestructura})}
\item contrato de fideicomiso; resolución de la Secretaría de Ambiente; certificado de aptitud ambiental municipal; sentencia completa de la Sala III \textit{(cap.\ \ref{cap:vaqueros})}
\item si las tareas del parte del 02/02/2026 fueron comunicadas al tribunal, si están contempladas en el Plan aprobado, y qué organismos provinciales aportaron y contrataron la maquinaria; y actas de activación del Comité Operativo de Emergencia durante el temporal de enero de 2026 \textit{(cap.\ \ref{cap:loteo})}
\item expediente 0090227-97704/2013-0: Estudio de Impacto Ambiental y Social de la Urbanización Cerro de Buena Vista, acta de la audiencia del 10/10/2014 y resolución final; \textbf{constancia de si existió estudio de impacto ambiental o audiencia pública para los loteos El Durazno y El Nogalar}, y en su caso su categorización conforme la Ley 7070; y estado actual del proyecto Cerro de Buena Vista ---la Ordenanza municipal 495/16 acredita la donación de calles y espacios verdes en 2016--- y plano de Inmuebles del expediente 18-28358 \textit{(cap.\ \ref{cap:loteo})}
\item registro municipal y provincial de alojamientos turísticos habilitados en La Caldera y Vaqueros, con plazas y ubicación; ordenanzas sobre alquiler temporario; serie completa de la Tasa Neta de Ocupación Turística departamental; y expediente del ``nuevo circuito turístico en la zona del dique'' \textit{(cap.\ \ref{cap:poblacion})}
\item \textbf{padrón inmobiliario del departamento} con valuaciones fiscales por parcela y su evolución 2010--2026; fecha del último revalúo general; y comparación entre valuación fiscal y precio de venta publicado de lotes en los loteos del corredor \textit{(cap.\ \ref{cap:hacienda})}
\item Decreto 4257/11 y convenio de 2010 con la Agencia Nacional de Seguridad Vial \textit{(el convenio del Decreto 1433/15 ya se obtuvo)}; rendición del destino a seguridad vial de lo percibido; monto anual efectivamente percibido por el municipio en concepto de formularios del Ce.N.A.T.; y estructura completa de ingresos municipales por origen --- tasas, coparticipación, convenios y contratos --- para los ejercicios 2014 a 2026 \textit{(cap.\ \ref{cap:hacienda})}
\item dotación de personal municipal en Obras Públicas, Ambiente y Tierra y Hábitat, con títulos y presupuesto, 2015 a la fecha; padrones inmobiliarios con titular de domicilio fiscal externo y su peso en la recaudación; recaudación efectiva del gravamen del art. 30 \textit{(cap.\ \ref{cap:hacienda})}
\item decretos del 9 de diciembre de 1983 que aceptan las renuncias de los presidentes de las Comisiones Municipales de La Caldera y Vaqueros, con los nombres de los renunciantes; decretos de designación de sus reemplazantes; texto del Decreto 877/82; ley que fijó la categoría de ambos municipios y su eventual modificación posterior; y verificación de si el Esteban Mogro\index{Mogro, Esteban} secretario administrativo y el juez de paz son la misma persona \textit{(cap.\ \ref{cap:politica})}
\item \textbf{acto de creación del Comité Metropolitano de Ordenamiento Territorial} del Área Metropolitana del Valle de Lerma, sus actas, el Registro Metropolitano de Urbanizaciones y las ordenanzas de los demás municipios ---La Caldera entre ellos--- que adoptaron el mismo manual de nuevas urbanizaciones \textit{(caps.\ \ref{cap:vaqueros} y \ref{cap:redes})}
\item \textbf{convocatoria a elecciones} para completar el mandato de intendente de Vaqueros conforme el art.\ 67 de la Ley 8126, o acto que la descarte \textit{(cap.\ \ref{cap:politica})}
\item acta de la sesión del Concejo Deliberante de Vaqueros del 01/03/2026; texto del fallo de la Corte de Justicia; Anexo I de la Ordenanza 507/16 (texto original del Código de Planeamiento Urbano Ambiental) y norma provincial que funda la jurisdicción de Vaqueros sobre Lesser \textit{(cap.\ \ref{cap:vaqueros})}
\item \textbf{decretos de designación} de los presidentes de la Comisión Municipal de La Caldera ---y, si los hubo, de sus intendentes--- entre 1983 y 1987 --- advertencia: los índices de decretos que el Boletín publica en línea traen sólo número, fecha y edición, \textbf{sin materia}, de modo que hay que ir por las ediciones o por el archivo del organismo (ver cap. \ref{cap:politica}); \textbf{actas de proclamación} de 1987, 1991, 1995, 1999, 2003 y 2007 --- no están en línea: el Tribunal Electoral publica desde 2009 y el Archivo Central Provincial no tiene documentos cargados, de modo que hay que pedirlas en mesa de entradas del Tribunal; \textbf{Decreto 1.288 del 18/05/2000} de la Secretaría General de la Gobernación, que aprueba un convenio entre ese organismo y el intendente de La Caldera --- primer año del primer mandato de Calabró ---, cuyo objeto no consta en el sumario del Boletín; verificación de los parentescos Calabró, Moreno Ovalle\index{Moreno Ovalle, Daniel Roberto} y Sumbay\index{Sumbay, Diego Ramiro} ante el Registro Civil; y composición completa de ambos Concejos Deliberantes por período \textit{(cap.\ \ref{cap:politica})}
\item \textbf{las capas georreferenciadas en formato digital} con las que el propio
Código declara haber confeccionado sus planos, en particular \textbf{la capa de la zona PU},
que es lo único que permite calcular el ancho de la franja: el articulado, la planilla del
Anexo VI y el plano del Anexo II ya se revisaron y \textbf{no contienen la cifra}. También los
ocho planos del atlas en soporte digital; ---las leyes 2131/47, 4365/70 y 4987/74 ya se obtuvieron y fijan el límite (capítulo \ref{cap:cot})---; \textbf{mensura del deslinde y amojonamiento} que el artículo 3 de la Ley 4987 encomendó a la Dirección General de Inmuebles en 1974; y el acto por el cual Lesser pasó al ejido de Vaqueros \textit{(cap.\ \ref{cap:cot})}
\item \textbf{si el límite que publica la capa «Límites Municipales» de IDESA coincide con el que el municipio reconoce} ---su dibujo sigue el Wierna, como la Ley 4987---, y \textbf{sobre qué ley se trazó}: es la comparación que permite verificar si el trazado publicado respeta la Ley 4987, en particular aguas arriba de la confluencia del Wierna con el río La Caldera \textit{(cap.\ \ref{cap:cot})}
\item verificar si el estudio de cota máxima edificable fue encargado y ejecutado \textit{(cap.\ \ref{cap:pdua})}
\item \textbf{texto de la Ordenanza Nº 643 de La Caldera tal como se sancionó}, y libro de actas del Concejo, para saber si aprobó sin cambios el proyecto de ordenanza del PDUA \textit{(cap.\ \ref{cap:pdua})}
\item \textbf{comprobantes de combustible de la emergencia de enero de 2026}, en particular el 0035-00006919, que figura en dos fechas, y filas 20 y 21 de la planilla \textit{(cap.\ \ref{cap:amparo})}
\item \textbf{actos municipales de habilitación de los loteos del sur del pueblo} ---El Nogalar, El Durazno, Potreros de La Caldera---, con las donaciones de remanentes y calles aceptadas, que el intendente atribuye en 2026 a gestiones anteriores: es la verificación que corresponde antes de dar por hechas esas atribuciones \textit{(cap.\ \ref{cap:amparo})}
\item \textbf{textos de las Ordenanzas 040/82 y 243/97} de La Caldera, que el proyecto de ordenanza del PDUA mantiene vigentes \textit{(cap.\ \ref{cap:pdua}; apéndice \ref{cap:normativa})}
\end{itemize}

\section{Boletín Oficial y archivo provincial}

\begin{itemize}[leftmargin=*]
\item \textbf{verificación sobre el facsímil de los hallazgos de 1947 y 1948}, que este libro publica leídos sólo en la capa de texto: el \textbf{Decreto 8176 del 5 de febrero de 1948} con los circuitos del departamento (Nº 3060 a 3065), el expediente \textbf{4887/A/1948} de reconocimiento de agua de Hermann Pfister\index{Pfister, Hermann} (Nº 3101, h.\ 4), el edicto de \textbf{«La Helvecia»} (Nº 3019 a 3023), el de \textbf{«El Angosto»} (Nº 3126 a 3132) y el decreto de \textbf{reelección} de Enrique Alejandro Pfister (Nº 3064, h.\ 2 y 6). \textbf{Ninguna hoja de esos dos años alcanzó el umbral de publicable del reconocimiento óptico}, y noventa de 1947 y ciento ochenta y siete de 1948 están marcadas para mirar enteras \textit{(caps.\ \ref{cap:siglo} y \ref{cap:fuentes})}
\item el \textbf{expediente 44-L del año 1929} completo, de la mina de plomo y plata «Caldera Nº 1, 2 y 3»: el padrón de 1933 la inscribe a nombre de Juan Larrán y la Resolución de Minas Nº 591 del 2 de abril de 1948 la adjudica a Enrique A.\ Zuccarino y señora. \textbf{El traspaso entre una cosa y la otra no se publicó}, y es lo que el expediente tiene que contar \textit{(cap.\ \ref{cap:siglo})}
\item el \textbf{informe posesorio de agua solicitado por don Luis Linares}, que el sumario de la edición 3104 de 1948 lista bajo \textbf{posesión treintañal} con el aviso Nº 3508, sin indicar departamento: es una vía distinta de la del reconocimiento ante la Administración General de Aguas, y el apellido es el mismo que el del peticionante de El Angosto \textit{(cap.\ \ref{cap:siglo})}
\item el \textbf{expediente y el plano de la «Fracción Entre Ríos»}, en Vaqueros, cuyo edicto de reconocimiento de agua publican tres ediciones de 1950: \textbf{cuándo se fraccionó la finca Entre Ríos, con qué plano y con qué aprobación}, y quién solicitó el agua \textit{(cap.\ \ref{cap:siglo})}
\item el \textbf{decreto 7321 del 13 de diciembre de 1947} completo, con su expediente: es el acto que \textbf{desafectó el Mercado Tipo I de La Caldera} invocando a «los propios habitantes» y que en el mismo texto \textbf{adjudicó el Mercado Tipo II de Pichanal} a Luis Moreno Díaz. Se pide, en particular, \textbf{la constancia de esa expresión de los habitantes} ---nota, petición, acta o expediente--- que el considerando invoca y no cita \textit{(cap.\ \ref{cap:siglo})}
\item la \textbf{ejecución de la partida del Ítem 6 de la Ley 834}: \$50.000 para defensas del río La Caldera y \$50.000 para el Wierna, ambos en primera etapa, presupuestados en 1947. Se pide saber \textbf{si se licitaron, se adjudicaron y se certificaron}, y con qué expediente. Y lo mismo de los \$370.000 del Inciso VI, fondos de la Ley 770, para la estación sanitaria de La Caldera y otras cinco localidades \textit{(cap.\ \ref{cap:siglo})}
\item las \textbf{fojas 12 a 18 del expediente 17847/1947}: el plan hidráulico orgánico de obras 1948--1950 de la Administración General de Aguas, aprobado por resolución 326 del acta 16 del 18 de julio de 1947 y por decreto 5466-E, con una inversión de \$35.100.000. \textbf{El decreto publica el monto y no el plan}, y de esas fojas depende saber si la defensa del río de la Caldera ---croquis en 1938, obra suspendida en 1940, ejecutada por vía administrativa en 1945--- tuvo continuidad \textit{(cap.\ \ref{cap:siglo})}
\item \textbf{con qué criterio la Administración General de Aguas distingue entre «temporal y permanente» y «temporal y eventual»}: los expedientes 4887/A/1948 (Pfister, La Helvecia) y 4498/47 (Linares, El Angosto) reconocen el mismo coeficiente de 0,525 litros por segundo por hectárea sobre el mismo río y les asignan caracteres distintos, sin que ningún acto lo explique \textit{(cap.\ \ref{cap:siglo})}
\item \textbf{en qué consta que los habitantes del departamento manifestaron que el mercado era ineficioso}: el decreto de diciembre de 1947 que desafecta el Mercado Tipo I de La Caldera lo invoca como motivo y no cita nota, petición, acta ni expediente. Y el \textbf{decreto 4945 del 1º de julio de 1947} completo, que lo había presupuestado \textit{(cap.\ \ref{cap:siglo})}
\item el \textbf{informe que la Administración de Vialidad debía elevar} sobre «\textit{en mérito a qué consideraciones ha sido incluido, fuera del plan, el camino denominado «La Caldera a Santa Rufina», en la red de caminos provinciales}», exigido por el decreto 5916-E del 22 de septiembre de 1947, y la \textbf{resolución Nº 4967} del H.\ Consejo de Vialidad que el Poder Ejecutivo exceptuó de la aprobación del Acta 188 \textit{(cap.\ \ref{cap:siglo})}
\item \textbf{cuál fue la observación de Contaduría General} sobre la licitación de las refecciones sanitarias de la escuela provincial de La Caldera, adjudicada a Conrado Marcuzzi por \$35.026,99: el decreto 12.581 del 19 de noviembre de 1948 la aprueba y el \textbf{12.744-E del 27 de noviembre insiste en Acuerdo de Ministros} contra esas observaciones, sin decir cuáles eran \textit{(cap.\ \ref{cap:siglo})}
\item el \textbf{legajo de Gerardo Guerrero} como receptor de rentas, y \textbf{por qué el Banco Provincial de Salta solicitó su designación} para el departamento (Decreto 8706-E, expte.\ 4567/F/1948). Y, si existió, \textbf{una localidad llamada «La Caldera» en el departamento Capital}: la Resolución 423-E de 1947 se la atribuye dos veces \textit{(cap.\ \ref{cap:siglo})}
\item \textbf{nota Nº 30 del 6 de septiembre de 1921 de la Junta Electoral de la Provincia} y la \textbf{resolución que declaró nula la elección del 7 de agosto} en el distrito de la Caldera y en otros once: el decreto 1828 las invoca y no las publica, y sin ellas no se sabe por qué se anuló ni si la elección de octubre llegó a hacerse \textit{(caps.\ \ref{cap:siglo} y \ref{cap:politica})}
\item las \textbf{seis ediciones de 1921 que el repositorio reemplaza por un marcador de ausencia} ---859, 872, 876, 878, 884 y 885--- y la hoja faltante de la 852 \textit{(apéndice \ref{cap:fuentes})}

\item \textbf{ediciones Nº 787 de 1919 y Nº 793, 822, 837 y 838 de 1920}, que faltan en el repositorio. La cadena de folios prueba que la \textbf{787}, del viernes 26 de diciembre de 1919, existe y tiene ocho hojas: 1919 cierra en el folio 1.072 y 1920 abre en el 1.081 \textit{(apéndice \ref{cap:fuentes})}
\item \textbf{acta de escrutinio y resolución de la Junta Escrutadora sobre la anulación de la urna de la mesa 4 del Colegio Electoral Nº 3 de Caldera} en la elección del 7 de marzo de 1920, que el decreto 772 invoca sin dar la causa; y \textbf{nota Nº 1.622 M.-19 de la Jefatura de Policía}, que funda el nombramiento del tercer sub-comisario de Mojotoro del mismo año \textit{(cap.\ \ref{cap:siglo})}

\item \textbf{edición Nº 737 del Boletín Oficial}, que falta en el repositorio y que es la única entre la extraordinaria del 31 de diciembre de 1918 y la Nº 738 del 17 de enero de 1919. \textbf{Y aclaración sobre la numeración de esos días}: la cadena de folios cierra el año 1918 en la página 587 y abre 1919 en la 595, de modo que a la 737 le corresponden siete hojas; pero entre las dos fechas hay \textbf{dos viernes} ---el 3 y el 10 de enero--- y una sola edición posible. O uno de los dos días no tuvo edición, o hay una edición más que no está en ninguno de los dos repositorios \textit{(cap.\ \ref{cap:siglo})}
\item las \textbf{38 hojas ausentes de 1919}, que la cadena de folios localiza una por una al final de veintiocho ediciones, y que son justamente las hojas donde el Boletín de la época imprime avisos y edictos \textit{(apéndice \ref{cap:fuentes})}

\item \textbf{expediente de la concesión de agua para regadío de la finca «Campo Alegre» de 1938}, con el informe de la Municipalidad de La Caldera y el dictamen del Fiscal de Gobierno: \textbf{si se concedió, con qué caudal y para cuánta superficie}. Es la segunda concesión de agua del departamento que el archivo registra, catorce años después de la de Urquiza \textit{(caps.\ \ref{cap:siglo} y \ref{cap:infraestructura})}
\item \textbf{criterio con que el Revalúo General de 1944 decidió qué partidas del departamento se publicaban}, y el padrón completo de La Caldera a esa fecha: el cuadro del Boletín Nº 2166 es un revalúo \textbf{provincial} con miles de partidas, y el tramo del departamento son las cuarenta y siete que el capítulo \ref{cap:fincas} transcribe. \textbf{La pregunta no es si están todas, sino qué criterio decidió cuáles entraban} \textit{(cap.\ \ref{cap:fincas})}
\item \textbf{acto por el que se dispuso la intervención de la comuna de La Caldera vigente en 1944}, y el nombre del Interventor: el decreto que aprueba el presupuesto de ese año lo da por existente ---«\textit{el señor Interventor de la comuna de la Caldera eleva}»--- y \textbf{ningún acto del año leído lo nombra ni dice cuándo se dispuso} \textit{(cap.\ \ref{cap:siglo})}. Se pide también \textbf{la edición 2.030 del Boletín}, del 26 de noviembre de 1943: es la única de esas semanas que no se obtuvo, y las otras cuatro, leídas al cierre, no traen el acto (apéndice \ref{cap:fuentes})
\item \textbf{acuerdo de la Municipalidad de La Caldera del 4 de junio de 1885} que concedió a Eustaquio Murúa treinta litros por segundo del río Wierna para la finca «Sausal o Curuzú», y \textbf{los demás derechos de agua que el municipio otorgó y que en 1942 la Provincia inscribió como propios} ---entre ellos los del río de las Nieves, reconocidos «desde fecha inmemorial»---. \textbf{Es el punto de partida de la voz municipal sobre el agua y el libro sólo tiene su cita de segunda mano} \textit{(caps.\ \ref{cap:planteo}, \ref{cap:siglo} y \ref{cap:infraestructura})}
\item \textbf{resultado del juicio de posesión treintañal del terreno de la comisaría de policía del pueblo de La Caldera}, deducido por el Gobierno de la Provincia en 1943, y \textbf{el informe de la Municipalidad sobre si afectaba propiedad municipal}: treinta años de posesión implican ocupación sin título desde 1913 o antes \textit{(cap.\ \ref{cap:siglo})}
\item \textbf{causa judicial por la muerte de Paulino Fernández}, agente de policía de la comisaría de La Caldera, que el expediente de pensión de 1941 atribuye a un disparo en el paraje La Calderilla: \textbf{qué resolvió la justicia}. El archivo leído trae el hecho sólo por la vía administrativa de la pensión, y sin la causa no se puede decir nada más que lo que ese trámite consigna \textit{(cap.\ \ref{cap:siglo})}
\item \textbf{expediente de adquisición de los terrenos para el camino Calderilla a río de las Pavas, tramo Desmonte--Campo Santo} (1941): \textbf{quiénes eran los propietarios, qué superficie se tomó y cuánto se pagó}. Es el primer acto en que una traza vial del departamento exige comprar el suelo, y el antecedente más antiguo de la relación entre camino y dominio \textit{(caps.\ \ref{cap:siglo} y \ref{cap:expropiacion})}
\item \textbf{expedientes de las tres defensas del río de la Caldera}: la de 1938, que quedó en croquis; el \textbf{reparo provisorio} de 1940 «en el lugar denominado la Calderilla»; y la de 1945 (expediente 18385, decreto 8549), que se llamó a licitación pública por \$2.515 y terminó ejecutándose \textbf{«por vía administrativa»}, sin que se publique certificado ni recepción. \textbf{Interesa saber si alguna de las tres llegó a construirse, dónde, y si queda rastro} ---es el antecedente directo del problema que el Plan de Manejo de 2024 cifra en el cincuenta y ocho por ciento de su inversión--- \textit{(caps.\ \ref{cap:siglo} y \ref{cap:amparo})}
\item \textbf{texto de la ordenanza municipal de La Caldera del 13 de octubre de 1914}, que en 1939 se invoca en un expediente provincial como norma vigente que «\textit{reglamenta la distribución de la totalidad}» del caudal, y \textbf{la resolución de la Provincia sobre esa oposición}: es la pieza que decide si el municipio conservaba o no potestad sobre el agua veinticinco años después de dictarla \textit{(caps.\ \ref{cap:siglo} y \ref{cap:infraestructura})}
\item \textbf{expediente de la obra «defensa río Caldera» de 1938}: los cómputos métricos, el presupuesto y \textbf{el croquis} que el decreto invoca, y sobre todo \textbf{si esos cincuenta metros de defensa sobre la margen izquierda llegaron a construirse y dónde están}. Es el antecedente más antiguo del problema que el Plan de Manejo de 2024 cifra en el cincuenta y ocho por ciento de su inversión \textit{(caps.\ \ref{cap:siglo} y \ref{cap:amparo})}
\item \textbf{balance del ejercicio 1937 de la Comisión Municipal de La Caldera} (expte.\ 1117 M) y el informe del Consejo de Educación que declara que no efectuó ingreso alguno al fondo propio de educación, con el \textbf{porcentaje} que el decreto de 1938 fija sobre lo recaudado. Ninguna de las tres piezas se publica \textit{(cap.\ \ref{cap:siglo})}
\item \textbf{expediente Nº 3463 D y los registros de la catastración de La Caldera que Adeodato Aybar declaró terminada en agosto de 1929}: superficie relevada, número de parcelas, planos y dónde quedaron, y el decreto que fijó finalmente su remuneración. \textbf{Es el antecedente más antiguo del catastro con el que este libro trabaja hoy} y de él sólo se publicó que estaba hecho \textit{(caps.\ \ref{cap:siglo} y \ref{cap:fincas})}
\item \textbf{actuaciones de la Municipalidad de Campo Santo bajo la Ley 11.078 de 1929}: qué trabajos hizo en la playa del río Mojotoro desde las juntas del Wierna con el río La Caldera, qué bocas-tomas prohibió y con qué criterio; y las \textbf{actas de la Comisión \textit{ad hoc} de octubre de 1931} que integraron los Comisionados Municipales de los dos departamentos sobre las aguas del Vaqueros y del Wierna \textit{(caps.\ \ref{cap:siglo} y \ref{cap:infraestructura})}
\item \textbf{constancia de la compra del terreno de la Casa Municipal de La Caldera} que autorizó la Ley 11.088 de 1929 por \$500: si se compró, dónde, y si es el mismo inmueble en que el municipio funciona hoy \textit{(cap.\ \ref{cap:siglo})}
\item \textbf{Registro Oficial de la Provincia del año 1928}, que resolvería el salto de numeración de decretos del 7097 al 8001 ---novecientos números en un solo día, con los dos tramos firmados por el mismo gobernador y el mismo ministro---: de esa lectura depende si los decretos 8023 y 9306 son de la serie corriente o de otra, y con ella la reconstrucción de la secuencia de mayo de 1928 \textit{(caps.\ \ref{cap:siglo} y \ref{cap:politica})}
\item \textbf{expediente 674-U} (renuncia de Francisco S.\ Urquiza a la Comisión Municipal) ---el 3898-C de la concesión figura completo, con su tramo de 1928, en un ítem propio de esta sección---; y el \textbf{fundamento del decreto 8023}, que deja sin efecto la concesión con la sola palabra «inconveniente» \textit{(cap.\ \ref{cap:politica})}
\item \textbf{memoria y rendición de cuentas de la Comisión Municipal de La Caldera por el ejercicio de 1926} (expediente 1099 M), y \textbf{presupuesto de gastos y cálculo de recursos de 1927} (expediente 1725 M), que los decretos 5024 y 5036 aprueban sin transcribir una sola cifra; y el \textbf{informe del Consejo General de Educación} sobre la \textbf{renta escolar} que le corresponde de la recaudación municipal, que los dos decretos invocan. \textbf{Es la primera rendición de cuentas municipal del departamento que el archivo registra, y de ella sólo se publicó que fue aprobada} \textit{(cap.\ \ref{cap:siglo})}
\item \textbf{resultado de los tres remates de la finca «Los Sauces» de 1927}, en el juicio sucesorio de Domingo Royo: si se vendió, a quién y en cuánto, después de que la base cayera de \$50.000 a \$14.000 en cinco semanas; y la \textbf{nómina de parajes del Colegio Electoral Nº 6 (Vaqueros)} de 1926 y 1927, que el decreto 3156 enumera a continuación de la del Nº 5 y que este libro no incorporó a la lámina \ref{fig:parajesnom} \textit{(cap.\ \ref{cap:siglo})}
\item \textbf{deslinde de la finca Las Nieves ejecutado por el agrimensor P.\ Arias en 1889}, cuyo vértice 49 usan como punto de arranque los dos edictos mineros de 1925 sobre esa finca (expedientes 1168-C y 1169-G). \textbf{Es el documento más antiguo que el archivo leído cita sobre el departamento}, treinta y seis años anterior al corpus, y no está \textit{(cap.\ \ref{cap:siglo})}
\item \textbf{acto por el que Francisco S.\ Urquiza fue nombrado subcomisario \textit{ad honorem} de Vaqueros}, que el decreto 2674 del 25 de junio de 1925 da por existente al reemplazarlo y que no aparece en ninguno de los quince tramos leídos: \textbf{interesa la fecha, porque de ella depende si el cargo policial es anterior o posterior a la concesión de agua de enero de 1924} \textit{(cap.\ \ref{cap:siglo})}
\item \textbf{estado de los cargos del departamento al cierre de 1925}: si Carlos Murúa conservó la Receptoría de Rentas y Recaudación de Impuestos al Consumo que le dio el decreto 2465 después de renunciar a la comisaría por el decreto 2962, sobre el que el acto calla; y el \textbf{presupuesto de reparaciones de la comisaría de La Caldera} por 3.302 m/l que aprueba el decreto 2152, cuyo detalle no se publica \textit{(cap.\ \ref{cap:siglo})}
\item \textbf{expediente Nº 5162 F}, donde corre el \textbf{Presupuesto de Gastos y Cálculo de Recursos de la Municipalidad de La Caldera para 1924} que el Decreto 1975 del 14 de noviembre aprueba sin publicar una sola cifra, junto con el \textbf{dictamen del Fiscal General} y el \textbf{informe del Consejo de Educación} que ese decreto invoca; y el \textbf{expediente Nº 4969 F}, con las «comunicaciones» que fundaron la designación de los cuatro miembros de la Comisión Municipal en abril de 1924 \textit{(cap.\ \ref{cap:siglo})}
\item \textbf{libros del Registro Civil de La Caldera que el inventario provincial del 14 de junio de 1935 no enumera}: ese cuadro registra para la oficina del departamento \textbf{un solo libro, el de matrimonios de 1913 a 1926}, mientras de las otras oficinas trae dos o tres. Interesa saber si existían los de nacimientos y defunciones, dónde están, y qué relación tiene esa única entrada con la diligencia de 1924 sobre «el destino que hayan tenido los libros» \textit{(cap.\ \ref{cap:siglo})}
\item \textbf{resultado de la diligencia ordenada por la resolución del Ministro de Gobierno del 1º de julio de 1924}, que autorizó a investigar «el destino que hayan tenido los libros» del Registro Civil del departamento: qué se encontró, y qué relación tiene con los cuatro encargados que ese Registro tuvo en cuatro meses de 1922. Ningún acto del año publica el resultado \textit{(cap.\ \ref{cap:siglo})}
\item \textbf{resultado de los dos remates de la finca «Calderilla» de 1924} ---el del 16 de abril sobre derechos y acciones, base \$533,33, y el del 11 de agosto sobre la mitad, base \$266,66---, cuyos avisos no coinciden en los titulares ni en el objeto ---el expediente 3898-C de la concesión de agua de 1924 figura en su propio ítem de esta sección--- \textit{(cap.\ \ref{cap:siglo})}
\item \textbf{expediente Nº 4541 F} y las presentaciones que instruye, que fundaron el \textbf{Decreto 1161 del 24 de septiembre de 1923 por el que la Provincia intervino la Municipalidad de La Caldera}; las \textbf{instrucciones} que el artículo 2º de ese decreto manda dar al Interventor por el Ministerio de Gobierno; y \textbf{los actos del Interventor Eduardo F.\ Harbin} y el acto que puso fin a la intervención, ninguno de los cuales se publicó en el año. \textbf{Es el acto de mayor alcance que el archivo temprano registra sobre el gobierno de este departamento, y lo único publicado es que ocurrió} \textit{(cap.\ \ref{cap:siglo})}
\item \textbf{informe de la Municipalidad de La Caldera en el expediente de la concesión de agua del río Wierna a Francisco S.\ Urquiza}, que el decreto 1.429 del 28 de enero de 1924 invoca como dictamen favorable de la Comisión Municipal: \textbf{interesa saber si lo firmó esa Comisión o el Interventor}, porque entre el pedido y el decreto la comuna fue intervenida. Y el \textbf{croquis del recorrido de la acequia} que el escrito de Urquiza dice acompañar y que el edicto no publica \textit{(cap.\ \ref{cap:siglo})}
\item \textbf{quince hojas de 1923 que la cadena de folios declara ausentes} del repositorio ---la única nombrable es el folio 3.334, entre las ediciones 943 y 944---, y \textbf{control de completitud de las veinte ediciones de la 971 a la 990}, donde la foliatura reinicia en cada número y 322 hojas quedan sin control \textit{(apéndice \ref{cap:fuentes})}
\item \textbf{ediciones del Boletín Oficial posteriores a la Nº 916 (21 de julio de 1922) y hasta el cierre de ese año} ---unas veintitrés---, que no están en el repositorio digital y que contendrían los actos del departamento del segundo semestre de 1922: entre ellos, \textbf{el resultado del remate de la finca Potrero del Gallinato fijado para el 28 de julio} y el acto que cubra la vacante que dejó en la Comisión Municipal la renuncia de Augusto Regis del 27 de junio \textit{(cap.\ \ref{cap:siglo})}
\item \textbf{presupuesto de gastos y cálculo de recursos de la Municipalidad de La Caldera para 1922}, sancionado por su Comisión Municipal en sesión del \textbf{28 de mayo de 1922} y aprobado por el Decreto 172 del 15 de julio, que lo aprueba sin transcribirlo; y el \textbf{dictamen del Fiscal General} que ese decreto invoca. Es el único presupuesto municipal del departamento que el archivo temprano menciona y no publica \textit{(cap.\ \ref{cap:siglo})}
\item \textbf{texto íntegro del decreto del Poder Ejecutivo de junio de 1922 que rectifica el artículo 2º del Decreto 510}, con su número propio, que no está impreso en la parte del acto que el repositorio trae: \textbf{qué se rectificaba} no consta, porque el decreto repite los mismos cinco nombres \textit{(cap.\ \ref{cap:siglo})}
\item \textbf{hojas faltantes dentro de las ediciones 890 a 916 de 1922} ---unas trece, que la cadena de folios impresos detecta y el conteo de archivos no---, y \textbf{la tapa de la edición 890}, cuya hoja 1 en el repositorio es una página interior \textit{(apéndice \ref{cap:fuentes})}
\item \textbf{edición Nº 703 del Boletín Oficial, del 17 de mayo de 1918}, que el repositorio digital reemplaza por un aviso de que no la encontró. Por su fecha contendría el \textbf{decreto de la Intervención Nacional del 11 de mayo de 1918} que el decreto N.\ 27 invoca como su fundamento y que no está en ninguna otra edición del año \textit{(cap.\ \ref{cap:siglo})}
\item \textbf{decreto N.\ 145 de junio de 1918}, citado por el N.\ 151 del 19 de junio y que no aparece publicado en un tramo donde no falta ninguna edición: o no se publicó, o se publicó fuera del Boletín \textit{(cap.\ \ref{cap:siglo})}

\item texto de la \textbf{Ley 5602/80 y su Decreto reglamentario 924/80} sobre clubes de campo, y de la \textbf{Resolución 25.819/97 de la Junta de Catastro} sobre urbanizaciones privadas: son los regímenes de los clubes de campo y de las urbanizaciones privadas que este libro documenta ---los loteos abiertos se aprobaron por la Ley 2.308---, y ninguno de los dos se obtuvo \textit{(cap.\ \ref{cap:loteo})}
\item acto del Poder Ejecutivo provincial que aprobó la Ordenanza 469/15 de ejidos urbanos de Vaqueros y Lesser, conforme el art.\ 98 de la Ley 1349 entonces vigente, y su inscripción en la Dirección General de Inmuebles. \textit{El anexo de la Ordenanza 437, el Código de Edificación 2026 y la propia 469/15 ya se obtuvieron del Digesto Jurídico Virtual de la Municipalidad de Vaqueros} \textit{(cap.\ \ref{cap:vaqueros})}
\item \textbf{sobre las leyes del departamento que el Digesto lista entre las no vigentes}, cuyos textos ya se obtuvieron: si el IPDUV construyó viviendas en el catastro 1229 de la Finca Vaqueros expropiado por la Ley 6334, o si operó la reversión; si se construyó la planta potabilizadora sobre los 3.710,08 m² del catastro 1280 (Ley 6464); qué ocurrió con la matrícula 75.465 del departamento Capital y con las defensas de la margen derecha del río Vaqueros que ordenó la Ley 6895; y destino actual del catastro 802 del barrio Jardín (Ley 7460) y de la matrícula 1846 (Ley 7699), con las obras hechas bajo el cargo de «desarrollo integral de la zona» \textit{(caps.\ \ref{cap:expropiacion}, \ref{cap:tierrafiscal} y \ref{cap:vaqueros})}
\item \textbf{vigencia de la Ley provincial 1269 (original 11078) de 1929} ---el Digesto de Diputados la lista entre las de objeto cumplido o no vigentes---, que autoriza a la Municipalidad de Campo Santo a reconcentrar el agua desde la confluencia del Wierna con el de La Caldera y a prohibir nuevas bocas-toma: desde cuándo y por qué acto dejó de regir, y si la Municipalidad de Campo Santo llegó a ejercer esas potestades \textit{(cap.\ \ref{cap:aguabaja})}
\item \textbf{la Separata del Boletín Oficial Nº 16.515, del 11 de noviembre de 2002}. \textbf{Su contenido ya se obtuvo} por el buscador de leyes del propio Boletín ---el Decreto Nº 1.989/02, su Anexo I de doce artículos y su Anexo II de definiciones se transcriben en el capítulo \ref{cap:ribera}---, de modo que lo que resta es el \textbf{ejemplar de la separata} para cotejar que el texto en línea sea completo, y \textbf{la constancia de su distribución}: la edición de ese día no trae el decreto, sino una línea que remite a ella
\item motivo de la ausencia de las ediciones posteriores a 2015 en el repositorio descargable, y existencia de un archivo alternativo de acceso masivo; y política de conservación y acceso digital del Boletín Oficial de la Provincia \textit{(cap.\ \ref{cap:opacidad})}
\item verificación del número de la ley que fija el régimen aplicado --- Ley 1336, renumerada 2614, y su modificatoria Ley 4272 de 1968, cuyo texto ya se obtuvo ---; expedientes de expropiación de la Ley 4486, si hubo demandas: montos depositados, valuaciones fiscales de 1972, sentencias y fechas de pago efectivo; planos de expropiación 00136, 00137, 00143, 00144 y 00148 ---el 00148 es probablemente el «plano 148» que el Decreto 1591/80 cita para el catastro 1575 (capítulo \ref{cap:tierrafiscal})---; y destino del remanente no inundado de cada una de las cinco propiedades \textit{(cap.\ \ref{cap:expropiacion})}
\item \textit{(El anexo del Decreto 1682/19 ya se obtuvo: está publicado en la edición web del Boletín.)} \textbf{Registro de la Unidad Coordinadora} (art.\ 11 del anexo) y resoluciones de autorización de preventa de los loteos del corredor, con sus garantías; \textbf{foja 59 del expediente de El Durazno}, con lo que certificó Recursos Hídricos; decretos de aprobación de los loteos del corredor dictados entre 1973 y 2019, y si alguno se acogió a la fianza del art.\ 13 del Decreto 1.410/73 \textit{(cap.\ \ref{cap:loteo})}
\item expedientes de los decretos de enero de 1910 y mayo de 1911, con la rendición de
cuentas de Daniel Linares\index{Linares, Daniel} ante la Contaduría, que diría qué reparos se construyeron y
dónde; aforos del río Wierna, si alguna vez se hicieron ---el expediente 3898-C, en su
propio ítem---; expedientes 1164 C, 1165 C y 1168 C del registro minero y límites de superficie
por solicitante en el Código de Minería de la época; liquidación de Levilly y Compañía y
resultado del remate de la fundición; y resultado del remate de Potrero del Gallinato de
1922 \textit{(cap.\ \ref{cap:siglo})}
\item \textbf{expediente 8045/1946} del Ministerio de Hacienda, Obras Públicas y Fomento ---informe de la Dirección General de Hidráulica de fs.\ 20 a 28 y cuadro de distribución de fs.\ 28--- y \textbf{Decreto 1654 H del 23/09/1946}, sobre el río Horcones, que el Decreto 1966 H invoca como precedente \textit{(cap.\ \ref{cap:redes})}
\item \textbf{actos provinciales de 1947 a 2012 que hayan dado, quitado o delegado a las comisiones municipales de La Caldera y Vaqueros competencia sobre el agua o el suelo de su jurisdicción}: el tramo del Boletín que este relevamiento no leyó de corrido \textit{(cap.\ \ref{cap:planteo})}
\item textos de las Leyes 7.322 y 6.994 y su reglamentación; acto por el cual se incorporó a La Caldera y a Vaqueros a la Región Metropolitana; y existencia de otras figuras metropolitanas provinciales --- residuos, agua, ordenamiento territorial --- que alcancen al corredor \textit{(cap.\ \ref{cap:redes})}
\item \textbf{expediente Nº 3898 letra C} completo, de la concesión de agua del río Wierna a Francisco S.\ Urquiza\index{Urquiza, Francisco S.}, en sus dos tramos: el de 1924, con el informe del Departamento de Obras Públicas y el de la Municipalidad; y el de 1928, con el reclamo de los vecinos y de la Comisión Municipal de Campo Santo, el informe de la Dirección de Obras Públicas que lo rechaza, los demás informes de ese año y quién los suscribió por la Comisión Municipal de La Caldera. Y el \textbf{Libro de Registro de Concesiones de Agua} de la Escribanía de Gobierno, donde el decreto de abril de 1928 ordena inscribir esa concesión. \textbf{Se unificaron aquí los cuatro pedidos de este expediente} que el apéndice tenía repartidos en la sección \textit{(caps.\ \ref{cap:siglo} y \ref{cap:politica})}
\item \textbf{Anexos de los Decretos 3069/14 y 463/15}: contratos de promoción turística con La Caldera S.A.\ y con Proyecto 1 S.A., con el detalle de exenciones, plazos y montos; expedientes 136-135.448/2013 y 136-34167/2010; \textbf{legajos societarios de ambas firmas} en la Dirección General de Sociedades; antecedentes dominiales de las matrículas 3.196 y 3.337; nómina completa de proyectos acogidos a la Ley 6064 en el departamento; \textbf{informes semestrales de la Dirección Provincial de Turismo} sobre ambos beneficiarios (art.\ 33 del Decreto 2448/85), con la nómina de socios que exige su art.\ 24 y el estado de la cancha de golf del comodato de 2015; y si la Municipalidad de La Caldera adhirió al régimen eximiendo tasas (art.\ 37) \textit{(cap.\ \ref{cap:poblacion})}
\item expediente C/24-02529/79 con la \textbf{evaluación de daños de la Dirección General
Agropecuaria}, que consignaría superficie tabacalera afectada y nómina de productores del
departamento en 1979; y serie de \textbf{declaraciones de emergencia agropecuaria} que hayan
alcanzado a La Caldera entre 1980 y 2018 \textit{(cap.\ \ref{cap:tierra})}
\item art.\ 28 inc.\ 12 y art.\ 33 de la Ley 8.171; \textbf{acto que trasladó la Secretaría a la Jefatura de Gabinete}, por la Ley 8.511 (B.O.\ del 26/11/2025): \textbf{manuales de misiones y funciones} de la Secretaría de Tierra y Bienes del Estado y de la Subsecretaría de Recursos Hídricos, ambas en esa jurisdicción según la Decisión Administrativa 569/25 (la Jefatura ya aprobó por las Resoluciones 26/26 y 47/26 los de las Secretarías de Relaciones Institucionales e Internacionales y del Interior); \textbf{foja 137 del expediente 18-37.348/2019-0}, que corresponde a la Secretaría de Tierras y Bienes del Estado; y ubicación actual de las competencias de la Unidad Coordinadora, dado que la Secretaría figura en 2025 como Secretaría de Tierra y Bienes del Estado bajo otra dependencia \textit{(cap.\ \ref{cap:loteo})}
\item \textbf{texto de la Ley 8274}, modificatoria de la Ley de Ministerios 8171, que el libro conoce sólo por los fundamentos del Decreto 198/24 \textit{(cap.\ \ref{cap:tierrafiscal})}
\item \textbf{antecedentes de la Destilería Chachapoyas} y de la jurisdicción policial propia que le reconoce el Decreto 5738-G de 1945 \textit{(cap.\ \ref{cap:infraestructura})}
\item \textbf{ordenanza municipal de Campo Santo del 18 de noviembre de 1866}, la ley que la sancionó en enero de 1867 y la \textbf{ley del 25 de septiembre de 1908} que dota de agua a General Güemes, que el fallo de 1910 cita sin transcribir \textit{(cap.\ \ref{cap:redes})}
\item \textbf{expedientes de cateo sobre la finca «El Cerro Nevado» de Garzón y Pintos anteriores a noviembre de 1916}: el de Casiano Hoyos «por sí y por don Juan Larrán»\index{Larrán, Juan} en el lugar «El Potrero», y el de Octavio Espeche y Odorico Díaz en la falda este del cerro, que el informe del Escribano de Minas cita sin acompañar; y el de José María Romero Escobar sobre el «Puesto de las Pirguas», para establecer si ese distrito estaba en el departamento \textit{(cap.\ \ref{cap:siglo})}
\item \textbf{número de Boletín y hoja de los treinta y un actos de dominio de 1924 a 1945} que la tabla del apéndice \ref{cap:dominio} registra sin cita de edición: el barrido de ese tramo no dejó asentada la edición fila por fila, y sin ella verificar una operación exige rehacer el año. Es el único defecto de localización que este libro se conoce a sí mismo \textit{(apéndice \ref{cap:dominio})}
\item \textbf{edición Nº 640 de 1917, ausente entera del repositorio, y las cuarenta hojas del año que faltan} respecto de las 329 que la paginación impresa declara \textit{(apéndice \ref{cap:fuentes})}
\item expediente del \textbf{deslinde y división de condominio de la finca «Potrero Gallinato»} de 1917 (Camilo Cáceres; juez de feria Augusto F.\ Torino; perito Herman Pfister\index{Pfister, Hermann}): \textbf{entre quiénes se dividió el condominio y con qué superficies}. Es el acto que explica cómo José Nogales pasó de lindero a dueño antes del remate de 1922 \textit{(cap.\ \ref{cap:siglo}, apéndice \ref{cap:dominio})}
\item \textbf{ediciones de 1915 que faltan en el repositorio digital}: la \textbf{Nº 581}, cuya descarga devuelve un archivo vacío, y las anteriores al 10 de febrero ---enero entero y la primera semana de febrero---, que son las únicas del año que este relevamiento no pudo leer \textit{(cap.\ \ref{cap:fuentes})}
\item \textbf{ejemplar impreso o escaneo a mayor resolución de la edición Nº 567 de 1915, hoja 2}: el primer dígito del tramo de padrón de la segunda mesa del departamento está empastado y admite leerse \textbf{216 o 316}. De cuál sea depende que el padrón de 1915 haya bajado o subido respecto de los 524 de 1912 \textit{(cap.\ \ref{cap:poblacion})}
\item \textbf{textos de los Decretos 450/18 y 1.045/96}: el primero encomienda al IPV la asignación de inmuebles fiscales; el segundo es el que invocan en su visto las Resoluciones 7, 8 y 9 de 2016 \textit{(apéndice \ref{cap:normativa})}
\end{itemize}

\section{Secretaría de Ambiente y Desarrollo Sustentable}

{\small\itshape Es la Autoridad de Aplicación del ordenamiento de bosques nativos. Estos siete pedidos son de una misma materia y van en una sola nota; el segundo tiene plazo vencido.\par}

\begin{itemize}[leftmargin=*,itemsep=0.35em]
\item \textbf{expediente 0090227-229822/2022-0}: resultado del sumario por el loteo sobre las matrículas 1555 y 161, auditorías ambientales presentadas por los loteos intimados y actuaciones posteriores a las actas de inspección 000001 y 000002 del 12/10/2022; y estado de la denuncia penal SCLC-283/2022 \textit{(cap.\ \ref{cap:amparo})}
\item \textbf{resultado de la identificación de desmontes ejecutados desde el primer OTBN y de su legalidad}, y \textbf{el plan de acción} del artículo 40 de la Ley 26.331, que el artículo 9º de la Ley 8483 mandaba concluir dentro de los 180 días de su promulgación ---\textbf{plazo vencido a fines de junio de 2025}--- \textit{(cap.\ \ref{cap:opacidad})}
\item \textbf{capa final del OTBN con los ajustes de la adenda provincial} \textit{(las capas del informe técnico ya se obtuvieron del repositorio y reproducen su Tabla 17)}, para confirmar que la versión aprobada coincide con la usada en el capítulo \textit{(cap.\ \ref{cap:opacidad})}
\item \textbf{sectores del departamento sin categoría en el OTBN} ---entre 22.400 y 26.800 hectáreas según el contorno, casi todas en el radio censal rural 0101 de Vaqueros y el resto en los radios 0201 y 0202; unas 19.700 son el catastro 102, fiscal y reclamado por la comunidad Kondorwaira---: si tienen bosque nativo y qué régimen los alcanza \textit{(las categorías del resto ya se calcularon)} \textit{(cap.\ \ref{cap:opacidad})}
\item \textbf{no objeción de la autoridad forestal} exigida por el artículo 20 de la Ley 7543 según texto de la Ley 8483, para las urbanizaciones del corredor posteriores a diciembre de 2024, y \textbf{nómina de las otorgadas y denegadas} en el departamento \textit{(cap.\ \ref{cap:opacidad})}
\item \textbf{capa digital del Mapa 11 del OTBN ---criterio 11--- para el departamento La Caldera} (el Anexo IX se publica como imagen), y el \textbf{Anexo VIII con el protocolo de consulta previa} y la composición de las Comisiones de Evaluación: \textbf{si el territorio relevado bajo la Resolución INAI 617/2013 está en esa capa}, todo proyecto de cambio de uso sobre él requiere consulta y visto bueno \textit{(cap.\ \ref{cap:opacidad})}
\item \textbf{cálculo del umbral máximo de deforestación de la cuenca Mojotoro--Lavayén} por el método del Número de Curva ---la Tabla 17 del informe ya la marca como excedida--- y \textbf{límite departamental oficial de la Provincia}, para confirmar el cruce que el capítulo hizo con el del IGN: tres cuartos del departamento en la cuenca \textit{(cap.\ \ref{cap:opacidad})}
\end{itemize}

\section{Organismos nacionales y censales}

\begin{itemize}[leftmargin=*]
\item \textbf{cuadros del Censo Nacional Agropecuario 2018 desagregados por departamento} ---escala de extensión, uso de la tierra y régimen de tenencia---, para completar la serie 2002, 2008 y 2018. Los totales, la escala de extensión y la superficie implantada de 2018 ya se obtuvieron del sistema de consultas (capítulo \ref{cap:tierra}); faltan el uso de la tierra y el régimen de tenencia \textit{(cap.\ \ref{cap:tierra})}
\item \textbf{procesamiento especial del CNA 2018 para el departamento La Caldera}: terrenos sin límites definidos por tipo de dominio y régimen de tenencia ---el cuadro 3.3 de los resultados definitivos sólo llega a nivel provincial---, para saber cuántas de las 61 explotaciones sin límites están en tierra fiscal o comunitaria \textit{(cap.\ \ref{cap:tierra})}
\item \textbf{qué pasó con las 75.870 hectáreas} que el CNA 2008 registraba como explotaciones agropecuarias con límites definidos y que en 2018 ya no se declaran: \textbf{si cambiaron de uso, de titular, o sólo dejaron de censarse} \textit{(cap.\ \ref{cap:tierra})}
\item \textbf{identificación de la explotación de más de diez mil hectáreas} que el CNA 2008 registra en el departamento: qué finca es y quién la declaraba \textit{(cap.\ \ref{cap:tierra})}
\item \textbf{actas de las mesas de concertación sectoriales del PDES 2030} en que participó la Intendencia de La Caldera, y \textbf{qué planteó} \textit{(cap.\ \ref{cap:trabajo})}
\item \textbf{Anexo I de la Resolución INAI Nº 617/2013}: cartografía georreferenciada del territorio relevado de la comunidad, el \textbf{acto que dio por concluido o prorrogó} ese relevamiento, y el \textbf{Dictamen Jurídico} de la carpeta técnica, que analiza la situación registral del territorio. La publicación del Anexo I no se encontró, y el convenio INAI--Salta de 2022 declara secreta la información del programa, de modo que la vía más directa es la comunidad. Interesa además \textbf{en qué modalidad se hizo el relevamiento}: el expediente provincial 0010007-67527/2019-28 lo llama «centralizado» en el informe del 10 de agosto de 2023 y «desde el área de descentralizada del INAI» en el dictamen del 4 de septiembre \textit{(cap.\ \ref{cap:resistencias})}; y \textbf{cuál es el expediente de personería de la Comunidad Kolla Urkuwasi}, que el mismo expediente numera 54-142017/2018-0 en un informe y 224-252732/18 en la invitación a la consulta
\item superficie implantada por cultivo del departamento La Caldera en 2002 y 2018, en particular tabaco, del sistema de consultas del Censo Nacional Agropecuario, para resolver la hipótesis energética; cuadro 5.2.1.4 de la Dirección General de Estadísticas provincial sobre superficie plantada de tabaco Virginia; y cartografía del OTBN sobre el departamento, para establecer qué proporción de las 87.357 hectáreas perdidas pasó a régimen de conservación \textit{(cap.\ \ref{cap:tierra})}
\item datos de hogares del Censo 2022 por gobierno local ---la separación de población y viviendas entre Vaqueros y La Caldera ya se hizo---; correlación con superficie efectivamente loteada en el período \textit{(cap.\ \ref{cap:poblacion})}
\item procesamiento del cuestionario básico del Censo 2010, área 66077, con variable condición de ocupación, para obtener la serie sin combinar fuentes \textit{(cap.\ \ref{cap:poblacion})}
\item misma relación hogares/viviendas sobre la base del Censo 2010; y desagregación de la vacancia por radio censal para aislar los loteos ---la fracción 66077-03 es El Gallinato y los loteos del sur están en la 02 (capítulo \ref{cap:ausencias})--- \textit{(cap.\ \ref{cap:poblacion})}
\item metodología de la clasificación M1/M2/M3; ley vigente que fija la categoría de ambos municipios; y si hubo proyecto de recategorización tras el Censo 2022. El cuadro de población por gobierno local del INDEC (cuadro 1.17) y las leyes de límites 2131/47 y 4987/74 ya se obtuvieron y no se piden \textit{(cap.\ \ref{cap:poblacion})}
\item procesamiento del Censo 2022, población que se reconoce indígena o descendiente, departamento La Caldera; consulta al Registro Nacional de Comunidades Indígenas y al registro provincial; informes de organizaciones de derechos humanos con trabajo territorial en el Valle de Lerma \textit{(cap.\ \ref{cap:resistencias})}
\item identificación de la explotación de 5.870 hectáreas que registra el CNA 2018: cuál de las fincas enumeradas por el artículo 6º del PDUA --- El Gallinato, San Alejo, Los Porongos, San Ignacio, Las Cañas o Potrero de Castilla --- corresponde a las 5.870 hectáreas, y quién es su titular \textit{(cap.\ \ref{cap:tierra})}
\item \textbf{ejecución efectiva del plan vial de quince años de la Ley 12.625} en el departamento y destino de la partida del Puente Río Caldera, en Vialidad Nacional; y \textbf{fecha de habilitación del primer puente de acceso al pueblo de La Caldera} \textit{(caps.\ \ref{cap:infraestructura} y \ref{cap:poblacion})}
\end{itemize}

\section{La vida cotidiana}

{\small\itshape Corresponden al capítulo \ref{cap:ausencias}. Ninguno se resuelve con un
expediente: son series que otros organismos llevan y que este relevamiento no pidió.\par}

\begin{itemize}
\item \textbf{nómina de entidades civiles con personería vigente} domiciliadas en el departamento, y \textbf{cuáles presentaron balances} en los últimos cinco años \textit{(cap.\ \ref{cap:ausencias})}
\item \textbf{padrón de establecimientos educativos} del departamento, con matrícula por año y por nivel, y \textbf{dónde cursan la secundaria} quienes terminan la primaria en los parajes \textit{(cap.\ \ref{cap:ausencias})}
\item \textbf{nómina de efectores de salud} del departamento, con dotación de personal, horarios y servicio de urgencias \textit{(cap.\ \ref{cap:ausencias})}
\item \textbf{padrón de comercios habilitados} de los municipios de La Caldera y Vaqueros \textit{(cap.\ \ref{cap:ausencias})}
\item \textbf{nómina de servicios de transporte de pasajeros} autorizados en el corredor, con recorridos, frecuencias y si alcanzan a los parajes \textit{(cap.\ \ref{cap:ausencias})}
\item \textbf{libros parroquiales} de bautismos, matrimonios y defunciones, y desde qué año se conservan: es la serie demográfica más antigua y continua del departamento \textit{(cap.\ \ref{cap:ausencias})}
\item \textbf{expediente, planos y pliego de la licitación del edificio para la Policía de La Caldera}, llamada el 15 de noviembre de 1909 por la Oficina de Inspección de Obras Públicas: si se adjudicó, a quién y qué se construyó \textit{(cap.\ \ref{cap:presencia})}
\item \textbf{prensa salteña de las temporadas de lluvias, de 1930 a 1960}, en la hemeroteca provincial: noticias de La Caldera incomunicada por crecidas y crónica social de veraneantes \textit{(cap.\ \ref{cap:poblacion})}
\item \textbf{procesamiento del censo 2010 y 2022 para el departamento}, en este orden de prioridad: \textbf{lugar de trabajo} de la población ocupada ---dentro o fuera del departamento---, que es la variable decisiva; \textbf{categoría ocupacional} ---asalariado, cuenta propia, patrón--- y nivel educativo alcanzado. \textbf{La rama de actividad ya se obtuvo} y está en el capítulo \ref{cap:trabajo}; lo que falta es el \textbf{lugar de trabajo}, que el procesamiento público no ofrece. Las tasas de 2010 y 2022 ya se obtuvieron por Redatam y están en los capítulos \ref{cap:poblacion} y \ref{cap:trabajo} \textit{(cap.\ \ref{cap:ausencias})}
\item personería jurídica, estatuto, fecha de constitución, actas y padrón de la \textbf{Asociación Civil de Pequeños Productores Agropecuarios del Departamento La Caldera}: cantidad de asociados, superficie declarada, producción y actividad reciente ---tiene asambleas publicadas hasta 2024---; título y antecedentes de ocupación del \textbf{Club de Regatas Güemes} sobre el Dique Campo Alegre; personería y actas de la Asociación de Vecinos de los Ríos Castellanos y Lesser; nómina completa de asociaciones civiles, clubes y centros vecinales con personería en el departamento; su actuación --- si la hubo --- durante los aluviones de 2018 y 2019 y en el amparo colectivo; identificación de sus firmantes; y relevamiento completo de asociaciones civiles con domicilio en el departamento, en el registro de Personas Jurídicas \textit{(cap.\ \ref{cap:resistencias})}
\item estadísticas del Centro de Mediación Comunitaria de La Caldera y del de Vaqueros: cantidad y materia de los casos, con especial atención a conflictos por tierra, agua y linderos; fecha de creación de ambos centros; y continuidad de los contratos después de 2015 \textit{(cap.\ \ref{cap:resistencias})}
\end{itemize}

\section{Obra pública, prestadoras y programas de financiamiento}

\begin{itemize}[leftmargin=*]
\item contrato de AGEMET S.R.L.\ y producto entregado para Vaqueros; expedientes 0110272-37093/2014-0 y 0110272-241674/2014-2; plano del proyecto de costanera contenido en el PDUA, y su relación con la avenida Costanera sur mencionada en el amparo; los tres municipios que precedieron a La Caldera en el programa, para establecer de cuál proviene la plantilla; términos de referencia del Programa de Mejora de la Gestión Municipal, préstamo BID 1855/OC-AR, y nómina de municipios salteños alcanzados; contrato del arquitecto Kuncar y producto entregado; y expediente 0110272-199.922/2013-0 \textit{(cap.\ \ref{cap:pdua})}
\item Expediente 136-221704/23: pliego, monto adjudicado, informe que funda el encuadre en el art.\ 15 inc.\ e), plazo de entrega, \textbf{fecha de firma del contrato} y \textbf{fuente de financiamiento} ---si se pagó con los fondos del CFI gestionados en el expediente 0090034-224086/2022---; antecedentes de YSATI Ingeniería S.R.L.\ y su relación con otras contrataciones hídricas provinciales; y \textbf{Expediente Judicial Nº 238670/20} ---probablemente el 23860/20 con un dígito de más---, para establecer su relación con el 780280 y en cuál de los dos interviene Jardín Celestial S.R.L.\ como tercero ---en el 780280 no figura entre los sujetos registrados--- \textit{(cap.\ \ref{cap:amparo})}
\item expediente 7807/15: adjudicación, monto, empresa, acta de recepción y \textbf{trazado efectivo de las redes colectoras}, para establecer qué proporción del ejido de Vaqueros quedó servida; \textbf{planos de mensura de servidumbre de paso y acueducto} presentados ante la Dirección General de Inmuebles, que identificarían nominalmente a los propietarios de los terrenos atravesados; y existencia de un proyecto equivalente para la localidad de La Caldera \textit{(cap.\ \ref{cap:redes})}
\item expediente 011-0068-271427/2014-0: ubicación catastral de las treinta viviendas de Vaqueros, estudio de suelos, empresa adjudicataria y acta de recepción; su relación con la línea de ribera del río Vaqueros determinada en 2014; y estado de esas viviendas después de los aluviones \textit{(cap.\ \ref{cap:tierrafiscal})}
\item lectura completa de la serie de prensa sobre el río Vaqueros 2020--2026 y verificación de sus fechas; expediente y proyecto ejecutivo del nuevo puente Vaqueros--Salta, con la línea de ribera considerada; convenios entre la Municipalidad de Salta, la de Vaqueros y la Provincia sobre encauzamiento; y ubicación del barrio 15 de Septiembre respecto de la línea de ribera del río Vaqueros \textit{(cap.\ \ref{cap:vaqueros})}
\item presentación de Naturgy NOA S.A.\ en el expediente EX-2025-112087756-APN-GD\#ENARGAS, y el Informe Técnico Intergerencial IF-2026-41299939 que la valora; consulta al ENARGAS sobre expedientes de expansión de red para el departamento La Caldera; estado de la red cloacal en ambos municipios; y cobertura efectiva de agua de red por radio censal, disponible en el Censo 2022 \textit{(cap.\ \ref{cap:redes})}
\item expediente 125-120.154/13 y resultado del Concurso de Precios 16/2014: adjudicataria, monto, acta de recepción y \textbf{trazado efectivo de las redes maestras}, para establecer si alcanzaron la zona alta de Vaqueros; y expedientes 5.836/14 y 5.438/14 de la empresa \textit{(cap.\ \ref{cap:redes})}
\item expediente 125-217.210/13: pliego, acta de recepción de la obra y estado actual de las defensas; expediente 19.199/08 del Juzgado de Minas sobre la cantera El Vaquero; y nómina de obras hídricas ejecutadas sobre el río La Caldera entre 2013 y 2018 \textit{(cap.\ \ref{cap:defensas})}
\item \textbf{Resolución 20/2022 de la Secretaría de Minería y Energía} ---ausente del Digesto Normativo Minero Ambiental de 2025 y del índice del Boletín Oficial; pedirla por nota a la Secretaría de Minería citando el expediente 99078-2021---, que aprobó el acuerdo con Aguas del Norte (expediente 99078-2021) sobre la zona de exclusión de 591 ha, con su plano, los planes de limpieza que la empresa debía informar y las actuaciones posteriores a las inspecciones que hallaron extractores sin permiso en esa zona, y plano completo de la empresa del 2 de diciembre de 2022 \textit{(cap.\ \ref{cap:amparo})}; padrón de conexiones categorizadas como Baldío en los ejidos de La Caldera y Vaqueros, serie mensual, ante COSAYSA o el Ente Regulador de los Servicios Públicos; y delimitación del área efectivamente servida por la empresa en el departamento \textit{(cap.\ \ref{cap:poblacion})}
\item expediente de la obra BAPIN 132888: contratista, certificados de avance, actas de paralización si las hubo, y estado a la fecha de publicación; y existencia de obras equivalentes para el resto del departamento \textit{(cap.\ \ref{cap:redes})}
\item \textbf{texto del artículo 15 de la Ley 8072}, con sus incisos, que la Provincia invoca para contratar con la consultora del Plan (inc.\ e) y con los municipios (inc.\ a); expediente 33-192.402/2014-0 y serie completa de contrataciones de Vialidad con la Municipalidad de La Caldera, incluidas las posteriores a 2017; certificados de obra y ubicación de los badenes construidos; y si alguno de esos trabajos recayó sobre los arroyos del amparo \textit{(cap.\ \ref{cap:hacienda})}
\item \textbf{expediente de la ampliación de red de agua del barrio Santiago Apóstol} anunciada en septiembre de 2026: cantidad de conexiones, monto, plazo y \textbf{qué barrios y parajes del departamento siguen sin red} \textit{(cap.\ \ref{cap:redes})}
\item acto de \textbf{adjudicación de la licitación del puente sobre el río Vaqueros y la
Circunvalación Este} (2022), y estado de ejecución de la obra; y nómina de contrataciones de
\textbf{DESA --- Defensa y Encauzamiento S.A.} en el departamento \textit{(cap.\ \ref{cap:redes})}
\item \textbf{textos íntegros de las Resoluciones 896/23, 314/24, 951/24, 300/26 y 479/26 de la Secretaría de Obras Públicas}: montos, adjudicatarios y fundamentos, ---la identidad de la obra rescindida en 2026 con la de la 951/24 ya está confirmada---, si alguna integraba el Plan de Manejo aprobado en febrero, y si el encauzamiento de la 300/26 es el proyecto que la reemplaza \textit{(cap.\ \ref{cap:defensas})}
\item expediente 366-16303/25 completo: pliego, presupuesto, selección del contratista, memoria técnica, estudio hidrológico de sustento y certificados de avance \textit{(cap.\ \ref{cap:amparo})}
\item \textbf{Estudio de Impacto Ambiental y Social del expediente 0110033-243466/2021-181}, obra «Pavimentación Ruta Nacional Nº 9, tramo Salta -- La Caldera», Dirección de Vialidad de Salta: \textbf{traza, secciones, inmuebles afectados, tratamiento de la línea de ribera} y cruces de cursos \textit{(cap.\ \ref{cap:vaqueros})}
\item \textbf{acta de la audiencia pública del 27 de noviembre de 2023} sobre esa obra, celebrada en el CIC de Vaqueros: \textbf{quiénes se presentaron como parte, qué observaciones se formularon y qué respuesta recibieron}; y si se presentó alguien domiciliado en La Caldera \textit{(cap.\ \ref{cap:vaqueros})}
\item \textbf{estado actual de la obra}: si se licitó, se adjudicó o se ejecuta, con expediente y monto; el Presupuesto 2026 asigna a la Ruta 9 entre Salta y La Caldera \$18.026 millones con aportes de Vialidad Nacional (capítulo \ref{cap:amparo}) \textit{(cap.\ \ref{cap:vaqueros})}
\item \textbf{ancho de zona de camino de la Ruta Nacional 9} en el tramo Salta--La Caldera, y \textbf{si está amojonado e inscripto} en las matrículas linderas \textit{(cap.\ \ref{cap:vaqueros})}
\item \textbf{nómina de emprendimientos privados sobre las laderas del corredor de la Ruta 9 que la nota de agosto de 2026 sobre las obras de Vialidad Nacional describe como carentes de autorización}; si el organismo sostiene esa afirmación, y el informe técnico en que se apoya; y \textbf{qué autorización vial les faltaría} ---permiso de acceso, ocupación de zona de camino, línea de edificación--- \textit{(cap.\ \ref{cap:vaqueros})}
\item \textbf{nómina de obras adjudicadas por la Municipalidad de la ciudad de Salta a INCOVI S.R.L.} entre 2017 y 2026, con montos y expedientes \textit{(cap.\ \ref{cap:infraestructura})}
\item \textbf{nómina de obras adjudicadas por la Municipalidad de la ciudad de Salta a \textsc{transporte hnos.\ mendoza s.r.l.}} entre 2022 y 2026, con montos y expedientes, y \textbf{si los pliegos especifican procedencia del material pétreo} \textit{(cap.\ \ref{cap:infraestructura})}
\item \textbf{extensión de la red de gas natural dentro del municipio de Vaqueros} ---casco, loteos de la periferia y clubes de campo---, con fecha de habilitación de cada tramo, ante Naturgy NOA S.A.\ y el ENARGAS; y \textbf{fecha de edición del atlas provincial del gas}; y la \textbf{leyenda de la capa} a la que corresponde el marcador del mapa interactivo sobre Vaqueros \textit{(cap.\ \ref{cap:redes})}
\item expediente de \textbf{autorización de desmonte} para la obra del acueducto Campo
Alegre y su \textbf{Plan de Forestación Compensatoria}, con superficie comprometida,
ejecutada y verificada. El Plan de Manejo fotografió los carteles: existen y tienen
responsable \textit{(cap.\ \ref{cap:expropiacion})}
\item \textbf{traza oficial del Acueducto Norte} y reproyección a coordenadas planas de los puntos del estudio de ruido de 2015 y de la Resolución 185/14 \textit{(cap.\ \ref{cap:expropiacion})}
\end{itemize}

\section{Juzgados civiles y comerciales, y Archivo General}

\begin{itemize}[leftmargin=*]
\item \textbf{texto del artículo 175 de la Constitución de la Provincia vigente en 1919} y del \textbf{artículo 16 de la Ley Orgánica de Municipalidades} que el decreto 33 de ese año da por derogado por aquél; y si la reforma constitucional que ese mismo decreto anuncia «\textit{para conseguir la creación de la Municipalidad electiva en todos los Departamentos}» llegó a intentarse. Es el origen documentado del umbral de cinco mil habitantes que la Ley 1349 de 1933 recogerá como techo de la tercera categoría \textit{(caps.\ \ref{cap:siglo}, \ref{cap:poblacion} y \ref{cap:politica})}

\item expediente del \textbf{juicio de posesión treintañal de 1943} sobre el predio de la Comisaría de Policía de La Caldera (Juzgado de 2ª Nominación en lo Civil, decreto del 18/02/1943): resultado, y si hubo oposiciones \textit{(cap.\ \ref{cap:tierrafiscal})}
\item \textbf{expediente 1744 del Archivo General}: la causa penal de 1928 contra Francisco S.\ Urquiza\index{Urquiza, Francisco S.} por «hurto y uso indebido de agua», con la denuncia de Juan A.\ Bianchi\index{Bianchi, Juan A.}, el auto de prisión preventiva y el sobreseimiento de 1930; y los autos de daños y perjuicios, con la pericia de fs.\ 63-65 y las declaraciones de fs.\ 31-44 tomadas por el Juez de Paz de La Caldera \textit{(cap.\ \ref{cap:politica})}
\item autos \textbf{«Francisco Moschetti y Compañía contra Rafael Delgado\index{Delgado, Rafael}»} ante el Juzgado de Comercio: dónde funcionaba el aserradero de La Caldera, desde cuándo, con qué monte se abastecía y quién lo compró en 1936 \textit{(cap.\ \ref{cap:tierra})}
\item autos \textbf{«Salvador Rosa contra Abel, Adolfo y Alberto Murúa»} ante el Juzgado de Comercio, secretaría del escribano Ferrari Sosa: quién compró la finca La Despensa en 1936, en cuánto, y por qué la base subió entre la primera y la segunda subasta \textit{(cap.\ \ref{cap:fincas})}
\item autos \textbf{«Embargo Preventivo Melchor Rocha vs.\ José González Yufra»} y resultado del remate de la estancia Potrero de Gallinato de octubre de 1930, con verificación de la base contra la imagen del Boletín \textit{(cap.\ \ref{cap:siglo})}
\item expediente 1274 del juicio ejecutivo por el que se remató en 1929 la casa quinta de Vaqueros: quién era el ejecutado y quién compró \textit{(cap.\ \ref{cap:poblacion})}
\item expediente 13906 de 1927, juicio ejecutivo del \textbf{Banco de la Nación Argentina}: quién era el ejecutado de la casa quinta de Lesser, y resultado del remate \textit{(cap.\ \ref{cap:poblacion})}
\item deslinde del \textbf{Puesto de las Garrapatas} y remates de la finca \textbf{Los Sauces}, ambos de 1927--1928 \textit{(hallazgos laterales del relevamiento del Boletín, no desarrollados en el cuerpo)}
\item expediente del juicio de \textbf{reposición de título de la finca Peñones o Los Peñones} (1926), promovido por el Dr.\ Francisco F.\ Sosa por María Isabel Gómez, y su resultado \textit{(cap.\ \ref{cap:politica})}
\item expediente del \textbf{remate judicial de la finca «Los Porongos hoy San Luis»} de 1916 (ejecución de Enrique Werner contra Humberto Graf von Garnier Turawa, juzgado de 1ª instancia en lo civil y comercial, juez Martín Gómez Rincón): \textbf{si la subasta del 15 de junio se realizó y quién compró}. Es el eslabón que falta entre el conde Garnier de Turawa y Daniel Linares\index{Linares, Daniel}, que figura como dueño en 1932, y el aviso declara 5.359 ha \textit{(cap.\ \ref{cap:fincas}, apéndice \ref{cap:dominio})}
\item expediente del \textbf{deslinde de la finca «Loma Bola»} promovido por Alberto Escudero en 1916, con el \textbf{plano de Herman Pfister}\index{Pfister, Hermann}, y el del \textbf{segundo deslinde de «Santa Mónica» y «Quitilipe»} del mismo año: superficie de cada una y, en el segundo, la \textbf{línea trazada por el ingeniero Pedro Cornejo el 9 de abril de 1889} que el edicto invoca como divisoria \textit{(apéndice \ref{cap:dominio})}
\item expediente del \textbf{juicio sucesorio de Asunción Espeche González de Wayar}, abierto ante el \textbf{Juzgado de Paz del departamento de Caldera} y declarado por edicto del 2 de septiembre de 1915 (B.O.\ Nº 568, h.\ 4): si el acervo comprende inmuebles del departamento y cuáles, porque es el único acto judicial de Caldera que deja el año y la tabla de dominio no registra ninguno \textit{(cap.\ \ref{cap:siglo}, apéndice \ref{cap:dominio})}
\item expediente del deslinde de 1932, con el \textbf{plano de Hermann Pfister\index{Pfister, Hermann}} y la
superficie de cada una de las cinco fincas; \textbf{cadena de titulares de ``El Durazno o
Peñones'' desde los Murúa hasta las matrículas 4366 y 6.558}; y padrón de propietarios
del departamento contemporáneo a ese deslinde \textit{(cap.\ \ref{cap:fincas})}
\end{itemize}

\section{Otros destinatarios}

{\small\itshape Legislatura, registro de cooperativas, Tribunal Electoral, IDESA, la propia comunidad Kondorwaira y pedidos con más de un destinatario.\par}

\begin{itemize}[leftmargin=*]
\item \textbf{concesión de energía hidráulica de Skru S.A.} al pie del dique Campo Alegre (expte.\ 0090034-179036/2018): resolución, permiso de uso de la presa, relación del catastro 1782 con el inmueble expropiado en 1972 y estado del proyecto \textit{(cap.\ \ref{cap:expropiacion})}
\item \textbf{resoluciones de otorgamiento o rechazo} de las concesiones de agua publicadas para el departamento entre 2004 y 2025 ---en particular las de los loteos Cerro de Buena Vista, La Misión Etapa III, Las Vertientes y Valle Hermoso---, y destino de la media hectárea eventual no asignada en el expediente 0090034-235.963/2022 (unas treinta gestiones; expedientes identificados en el cap.\ \ref{cap:aguabaja}), con caudal, uso y carácter de cada una \textit{(cap.\ \ref{cap:aguabaja})}
\item \textbf{concesión de agua del Club de Campo Las Vertientes--Eco Pueblo}: resolución recaída en el expediente 0090034-263317/2015 (600,48 m³/día del río Wierna para 630 familias), caudal otorgado y relación entre el catastro 4095 y la matrícula 3989 \textit{(cap.\ \ref{cap:redes})}
\item \textbf{serie de líneas de ribera del arroyo Chaile o Teodolinda (2016--2024)}: expedientes completos, solicitantes, \textbf{titulares de los catastros 1000, 1095, 1750, 2318, 2320, 3604 y 5928 y de las matrículas 1534, 2340 y 2821}, ubicación del catastro 3604 (los demás se ubicaron en la capa catastral); coordenadas y zona inundable de la Resolución 75/2020; relación de estas determinaciones con loteos o construcciones en Vaqueros; y corrección de la coordenada publicada con un dígito de menos en la Resolución 205/17 \textit{(cap.\ \ref{cap:ribera})}
\item \textbf{Corredor Intermunicipal para el Desarrollo Sustentable (CorInDeS)}: protocolo adicional con su estatuto, ordenanzas de aprobación de La Caldera y Vaqueros, actas y resoluciones de la Asamblea de Intendentes desde 2009, a las municipalidades firmantes y a la representación provincial ante organismos internacionales \textit{(cap.\ \ref{cap:redes})}
\item \textbf{superficie legal del departamento La Caldera y de cada uno de sus dos municipios}, con el plano y las leyes de límites con que se calcula, a la Dirección General de Inmuebles \textit{(cap.\ \ref{cap:fincas})}
\item \textbf{cuadro 1.x definitivo del Censo 2022 de la Ciudad de Buenos Aires}, única jurisdicción no cotejada para el cuarto puesto nacional de crecimiento \textit{(cap.\ \ref{cap:poblacion})}
\item \textbf{constancia del cese de Héctor Miguel Calabró como Secretario de Asuntos Municipales} en diciembre de 2019: la búsqueda de decretos por su nombre no devuelve un acto individual, de modo que el cese pudo operar con el cambio de gestión; los de designación ya se obtuvieron (Decretos 89/15 y 1634/17) \textit{(cap.\ \ref{cap:loteo})}
\item \textbf{planilla original de electores de 2017} por municipio, para verificar la cifra de 12.299, idéntica a la población censal de 2022 \textit{(cap.\ \ref{cap:politica})}
\item \textbf{texto original del manifiesto vecinal de octubre de 2018}, con su título y la nómina de firmantes, a los vecinos autoconvocados que lo difundieron o a los medios que lo publicaron \textit{(cap.\ \ref{cap:resistencias})}
\item \textbf{nómina de asociados de la cooperativa matrícula 22.359 a octubre de 2018}, en el registro de cooperativas. Es el único documento que permite verificar o descartar la imputación personal que contiene la gacetilla vecinal de ese año, y \textbf{no debe publicarse nada al respecto sin haber cursado antes la carta de derecho a réplica}. En el mismo registro, su \textbf{estatuto y denominación social vigente}: si «Cooperativa de Provisión Minera La Mesa Redonda Ltda.» y «Cooperativa de Provisión de Áridos La Mesa Redonda Ltda.» son la misma entidad, y desde cuándo \textit{(cap.\ \ref{cap:resistencias})}
\item \textbf{convenio entre la autoridad minera provincial y la Municipalidad de La Caldera} por el que se delega el control de la extracción de áridos y el cobro del canon: texto, fecha, vigencia y \textbf{actas de inspección labradas en su cumplimiento} \textit{(cap.\ \ref{cap:resistencias})}
\item \textbf{servicios de datos de IDESA} ---direcciones \textsc{wms}, \textsc{wfs} o equivalentes--- para la capa «Catastro Urbano y Rural», y \textbf{constancia de si los atributos parcelarios son públicos}; en su caso, el procedimiento para consultar el parcelario por matrícula; y \textbf{por qué el catastro de la capital se publica con etiquetas y el de los departamentos no} \textit{(cap.\ \ref{cap:opacidad})}
\item el \textbf{``Plan de Mínimas''} elaborado por la Subsecretaría de Recursos Hídricos «independientemente del Plan de Manejo», mencionado dos veces en la audiencia del 12 de febrero de 2026: qué obras prevé, cuándo se elaboró y qué se ejecutó de él; el capítulo documenta el programa desde 2010 y el Acta Acuerdo de 2016, pero no la versión que se invocó en 2026 \textit{(cap.\ \ref{cap:amparo})}
\item \textbf{registros de la estación de la UCASAL} que monitorea la tendencia hidrogeológica de los cauces, y \textbf{planificación de obras prioritarias} que la Provincia debía presentar dentro de los cuarenta y cinco días de la aprobación del Plan \textit{(cap.\ \ref{cap:amparo})}
\item \textbf{carpeta técnica de relevamiento territorial de la Asociación Comunitaria Kondorwaira}, con su cartografía georreferenciada y su dictamen jurídico sobre la situación registral. \textbf{Fue entregada a la comunidad en acto público el 25 de noviembre de 2016}, de modo que la vía más directa es la propia comunidad; subsidiariamente, el INAI o la Subsecretaría de Regularización Territorial de Salta. Con esa carta puede establecerse si alguno de los loteos documentados se superpone con territorio de ocupación tradicional: \textbf{es la pregunta abierta más importante de este trabajo}. El cruce a nivel de catastro ya se hizo: ninguno de los loteos documentados toca el catastro 102, fiscal, que Pérez Domínguez y Belmonte (2024) identifican como el territorio reclamado; lo que falta es el perímetro fino de la carpeta \textit{(cap.\ \ref{cap:resistencias})}
\item identificación registral de \textbf{Bernardo Austerlitz\index{Austerlitz, Bernardo}}: si el propietario de San Alejo y Santa Rufina de 1925 y el senador por Caldera de 1927 son la misma persona, y cadena de titulares de esa finca desde entonces \textit{(caps.\ \ref{cap:siglo} y \ref{cap:politica})}
\item listado de áreas protegidas del Sistema Provincial y verificación de si alguna comprende superficie del departamento La Caldera; convenios del art.\ 26 sobre transferencia de terrenos fiscales a municipios con fines de conservación; adhesiones de propietarios privados del art.\ 47 inc.\ 2 en el departamento; y anotaciones registrales del art.\ 57 \textit{(cap.\ \ref{cap:tierra})}
\item texto de la Ley provincial 4.736 y su vigencia; actuaciones del expediente 90-31.780/2023 en la Comisión de Minería, Recursos Naturales y Medio Ambiente, si las hubo; cartografía de la subcuenca del río Porongal y su relación con la finca Los Porongos; declaraciones juradas patrimoniales del período 2021--2025; registro de excusaciones; y \textbf{pedido de derecho a réplica} \textit{(cap.\ \ref{cap:politica})}
\item \textbf{los contratos ratificados por los Decretos 3069/14 y 463/15}, con el detalle
de exenciones, plazos, montos y contraprestaciones, y la nómina completa de contratos de
promoción turística celebrados en el departamento; \textbf{procesamiento del censo por categoría ocupacional} para el departamento ---la rama y la condición de actividad ya se obtuvieron---, y \textbf{cantidad de puestos de
trabajo declarados} por las concesiones de áridos y por los alojamientos habilitados \textit{(cap.\ \ref{cap:trabajo})}
\item \textbf{dónde llega efectivamente la red de agua y de cloaca} en los ejidos de La
Caldera y Vaqueros, y desde cuándo: las proporciones de hogares sin servicio ya se obtuvieron del
Censo 2022, de modo que lo que falta es la traza y la serie de altas, ante la prestadora o el
ente regulador; \textbf{caudal
efectivamente captado} por el embalse Campo Alegre y por el Acueducto Norte, y su destino por
localidad; \textbf{nómina de concesiones de agua pública vigentes} sobre los cursos del
departamento, con caudal, titular y carácter permanente o eventual; el \textbf{aforo ``De
Partida'' del artículo 258} para las fuentes de esta cuenca, con su fecha y sus renovaciones
decenales; y el \textbf{inventario de aguas del artículo 161}, con la constancia de sus
actualizaciones anuales \textit{(cap.\ \ref{cap:aguabaja})}
\item \textbf{votos nulos y participación oficial} por mesa del departamento La Caldera en las elecciones de 2025 ---el padrón por mesa ya se obtuvo---; \textbf{altas por cambio de domicilio} registradas en el padrón del departamento desde 2011, por año; y, si la ley lo permite, la proporción de electores cuyo domicilio documental está en el departamento pero figuran censados en otro \textit{(cap.\ \ref{cap:politica})}
\item \textbf{registro de las conversaciones con dos vecinos sobre el veraneo}: fecha, si lo relatado es recuerdo propio o transmitido y por quién, y consentimiento para citarlas sin nombre, a los propios vecinos \textit{(cap.\ \ref{cap:poblacion})}
\end{itemize}
